% Self-contained source generated from article.tex and its distributed sections.
\documentclass[11pt,a4paper]{article}
\usepackage[T1]{fontenc}
\usepackage{lmodern}
\usepackage[margin=25mm,headheight=14pt]{geometry}
\usepackage{amsmath,amssymb,amsthm,mathtools,mathrsfs}
\usepackage{booktabs,longtable,array,enumitem,microtype,xcolor,xurl}
\usepackage{fancyhdr}
\definecolor{inkblue}{RGB}{26,57,89}
\usepackage[colorlinks=true,linkcolor=inkblue,citecolor=inkblue,urlcolor=inkblue]{hyperref}
\usepackage{bookmark}
\emergencystretch=3em
\allowdisplaybreaks[1]
\setcounter{tocdepth}{2}
\numberwithin{equation}{section}
\theoremstyle{plain}
\newtheorem{theorem}{Theorem}[section]
\newtheorem{lemma}[theorem]{Lemma}
\newtheorem{proposition}[theorem]{Proposition}
\newtheorem{corollary}[theorem]{Corollary}
\newtheorem{conjecture}[theorem]{Conjecture}
\theoremstyle{definition}
\newtheorem{definition}[theorem]{Definition}
\theoremstyle{remark}
\newtheorem{remark}[theorem]{Remark}
\newtheorem{example}[theorem]{Example}
\DeclareMathOperator{\Li}{Li}
\DeclareMathOperator{\Res}{Res}
\DeclareMathOperator{\rank}{rank}
\DeclareMathOperator{\sech}{sech}
\newcommand{\src}[1]{\nolinkurl{#1}}
\newcommand{\sourcepin}{\texttt{73f18ef013d5b74eafc883efb2c15210642df12a}}
\newcommand{\shortpin}{\texttt{73f18ef}}
\pagestyle{fancy}
\fancyhf{}
\fancyhead[L]{\small Harmonic tails and spectral Stieltjes calculus}
\fancyhead[R]{\small\thepage}
\renewcommand{\headrulewidth}{0.3pt}
\title{\textbf{Harmonic-Tail Reconstruction,\\
Quartic Trigamma Moments,\\and Spectral Stieltjes Calculus}\\[8pt]
\large Exact continuations of the ProveIt polylogarithm programme}
\author{Research continuation prepared for Vladimir Reshetnikov}
\date{11 October 2026}
\hypersetup{pdftitle={Harmonic-Tail Reconstruction, Quartic Trigamma Moments, and Spectral Stieltjes Calculus},pdfauthor={Research continuation prepared for Vladimir Reshetnikov},pdfsubject={Exact identities, harmonic reconstruction, polylogarithms and spectral Stieltjes derivatives}}
\begin{document}
\maketitle
\begin{abstract}
We continue several specific problems in the ProveIt manuscript and its
incoming reports.  First, we complete optimal-prefix reconstruction for
every unequal even pair with smaller tail order eight.  A structural
factorization identifies the highest Bernoulli-degree determinant with a
strictly positive zeta moment determinant; an exact finite certificate
closes the residue class left open in the preceding report.  The same
factorization proves effective eventual reconstruction for every fixed
smaller even order, reducing each fixed-order conjecture to finitely many
explicit rational matrices.

Second, we evaluate every regular centered quartic trigamma moment after
prescribed polynomial subtraction.  The values lie in six specified
formal coordinates, and a finite Bernoulli transform, an anchored analytic
generator, and a finite polylogarithm construction of depth at most three
are given.  Hoffman's classical unweighted quartic evaluation supplies
one starting constant and is explicitly credited.

Third, we evaluate the second spectral factorization defect for any
number of quadratic masses by ordinary logarithmic sums and an elementary
one-dimensional resolvent integral, including every confluent coefficient.
Finite formulas at all higher spectral derivatives retain the otherwise
lost Hurwitz jets.  A one-sided counterpart identifies all generalized
Stieltjes corrections, treats the source's conjugate-linear factorization
problem, and gives an anchored dilogarithm--Gamma primitive.
We also correct a research-status overlap concerning the already
established entire multiple-harmonic interpolant.  Proofs, exact finite
certificates, independent numerical diagnostics, provenance, and a
concrete further-research agenda accompany the article.
\end{abstract}
\begin{center}
\small Source snapshot: \sourcepin.\\
All claims of continuation are relative to this inspected snapshot;
the article is a research draft with explicit proofs and replay artifacts.
\end{center}
\clearpage
\begingroup\small
\tableofcontents
\endgroup
\clearpage

% BEGIN EMBEDDED sections/01_scope.tex
\section{Scope, exact targets, and conventions}
\label{sec:scope}

The useful unit of progress in this programme is a proved identity or a
precisely delimited structural theorem.  Several nearby special-value
problems already have substantial proofs, and several plausible-looking
new directions are already present in earlier reports.  Accordingly,
the results below are attached to specific source questions before their
proofs are developed.  Bounds appear only inside the reconstruction and
convergence proofs; they are not independent objectives of the article.

The principal new conclusions are the complete order-eight theorem,
the effective fixed-order theorem, the all-moment quartic evaluation,
and the higher spectral factorization calculus.  These conclusions have
different logical content.  A nonzero rational determinant proves a
coefficient-rank statement; it does not prove independence of the real
constants later substituted into the coordinates.  A normally convergent
subtracted sum is an ordinary identity; the divergent raw sum preceding
that subtraction is not assigned an unspecified value.

\subsection{The concrete questions being continued}

\begin{longtable}{@{}>{\raggedright\arraybackslash}p{0.29\textwidth}>{\raggedright\arraybackslash}p{0.35\textwidth}>{\raggedright\arraybackslash}p{0.30\textwidth}@{}}
\toprule
Source target & Result here & Boundary retained\\
\midrule
\endhead
Reconstruction report, Q1: the omitted class
$b\equiv1\pmod{16}$ for $(8,2b)$ \cite{Reconstruction}
& Theorem~\ref{r8c:main} covers every $b\ge5$.
& The unrelated Gaussian period conjectures are unaffected.\\[5pt]
Reconstruction report, Q2: general unequal even orders
& Theorem~\ref{ert:main} gives a terminating threshold algorithm for
each fixed smaller order.
& The remaining finite set is not checked for every smaller order.\\[5pt]
Reconstruction report, Q5: a quartic formula, fixed space, and generator
& Theorems~\ref{qtm:main} and \ref{qtm:generator} give the complete
regular parity family and six formal coordinates.
& Fifth and higher powers require additional analysis.\\[5pt]
Quadratic-spectrum report, Q7: several factors beyond the first derivative
\cite{Quadratic}
& Theorems~\ref{msj:thm:second}, \ref{msj:thm:quadrature}, and
\ref{msj:thm:confluent} give the cross term and its full confluent series.
& A finite classical-polylogarithm reduction of the integral at arbitrary
parameters is not asserted.\\[5pt]
Quadratic-spectrum report, Q2: higher conjugate-linear defects
& Corollary~\ref{cor:linear-jets} gives every spectral order, with the
correct scale and all Stieltjes terms.
& The first-derivative parity criterion does not eliminate higher
nonlocal terms.\\[5pt]
Same-direction report, Q2: entire higher-depth connected interpolation
\cite{SameDirection}
& Section~\ref{sec:audit} identifies the existing all-depth interpolant
and geometric kernel exactly.
& This is an attribution and research-status correction, not a new
entire-family theorem.\\
\bottomrule
\end{longtable}

Two additional constructions go beyond those literal questions.
Theorem~\ref{thm:alljets} isolates the Hurwitz-pole contribution at every
derivative order for arbitrary polynomial weights and several quadratic
factors.  Theorem~\ref{thm:primitive} integrates a nonlinear second-jet
identity explicitly in the shift, using dilogarithms and a Gamma ratio.
Both have their normalization fixed in the statement.

\subsection{Source pin and attribution}

The repository was inspected at commit \sourcepin, including the
canonical manuscript status files and the relevant harmonic, Gaussian,
Hurwitz, and spectral branches.  All six archives currently in
\src{docs/incoming} at that pin were extracted and triaged.  The detailed
proof development uses the reconstruction, quadratic-spectrum, and
same-direction reports; the Beta-crossing, Gamma-continuation, and
shift-and-twist reports were also considered when selecting targets.
The exact source paths and SHA-256 digests are recorded in
\src{provenance.json}.  This is not an assertion that every theorem in
the four-volume manuscript has been audited.

Classical ingredients include Bernoulli--Hurwitz asymptotics, Euler's
Bernoulli formula, polygamma partial fractions, Gamma reflection, and
normalized Mellin continuation \cite{DLMF}.  Hoffman's quartic tail
sum \cite{Hoffman}, the Ihara--Nakamura--Yamamoto interpolant
\cite{INY}, and the first-derivative multiplicative-anomaly framework
\cite{EVZ,PolynomialAnomaly} are attributed where used.  Inherited
determinant reductions and cubic formulas are identified explicitly.
No claim of exhaustive historical priority is inferred from absence in
the inspected repository.

\subsection{Common conventions}

We use
\[
 H_n^{(s)}=\sum_{k=1}^n k^{-s},\qquad H_0^{(s)}=0,\qquad
 A_s(n)=\zeta(s,n+1),\qquad x_n=n+\tfrac12
\]
in the centered harmonic sections.  For an integer $s\ge2$,
$A_s(n)=\zeta(s)-H_n^{(s)}$.  The polygamma convention is
\[
 \psi^{(m)}(z)=(-1)^{m+1}m!\,\zeta(m+1,z),\qquad m\ge1.
\]
In particular $A_2(n)=\psi'(n+1)$ and
$A_{2k}(n)=\psi^{(2k-1)}(n+1)/(2k-1)!$.

Bernoulli polynomials are normalized by
\[
 \frac{te^{xt}}{e^t-1}=\sum_{m\ge0}B_m(x)\frac{t^m}{m!},
 \qquad B_m=B_m(0),\qquad B_1=-\tfrac12.
\]
The crucial centered zero is
\[
 \zeta(-2j,1/2)=0\qquad(j\ge0).
\]
It makes the specified polynomial subtractions have zero continued
sum.  A shifted analogue must restore its nonzero Hurwitz values;
this midpoint normalization is not silently transported to other shifts.

Superscripts on $\zeta$ and on spectral families denote derivatives
in their first, spectral variable unless stated otherwise.  We use
\[
 \zeta(1+u,a)=\frac1u+
       \sum_{m\ge0}\frac{(-1)^m}{m!}\gamma_m(a)u^m,
 \qquad \gamma_0(a)=-\psi(a).
\]
Thus the order derivative of $\zeta$ and the parameter derivative of a
Stieltjes constant are different operations.  All exchanges of spectral
differentiation and infinite sums below follow a locally uniform
subtraction argument, rather than formal differentiation at a pole.

The one-variable multiple polylogarithm uses the outer-index argument:
\[
 \Li_{s_1,\ldots,s_d}(z)=
 \sum_{n_1>\cdots>n_d\ge1}
       \frac{z^{n_1}}{n_1^{s_1}\cdots n_d^{s_d}}.
\]
At $z=1$ with admissible positive integer indices this is the
descending-index multiple zeta value.  In the finite generator formulas
we first take $z=e^{-t}$ with real $t>0$; the anchored integral formulas
specify the holomorphic continuation through the cancellations at $t=0$.

\subsection{Why the proofs use subtractions and leading determinants}

The harmonic and spectral problems share a finite algebraic step but have
different analytic mechanisms.  For reconstruction, the source matrix is
indefinite.  Its highest degree in a fixed number of Bernoulli variables
nevertheless factors into a positive discrete moment determinant.  The
factorial size of the Bernoulli variables makes that highest-degree part
dominate eventually.  The exact exceptional arithmetic remains visible
and is handled separately for order eight.

For quartic moments, summation by parts reduces the problem from the
fourth tail power to the inherited cubic coordinates.  A specific
boundary cancellation removes the remaining cubic Bernoulli convolution.
The surviving lower-power terms explain all exceptional low indices;
they cannot be discarded by parity alone.

For spectral derivatives, the common-scale difference vanishes to order
two in the spectral parameter before summation.  A resonant Hurwitz
pole can lower that order by one.  Its residue, regular Stieltjes
coefficients, and higher coefficient of the vanishing factor all matter.
The finite subtraction formulas retain those contributions while leaving
an ordinary summable remainder.  This is the source of the extra local
term in \eqref{one:x5second} and of the higher bilateral correction in
\eqref{msj:eq:third-contact}.

% END EMBEDDED sections/01_scope.tex


% BEGIN EMBEDDED sections/02_rank_eight.tex
% Research fragment for the parent article. All labels use r8c:.
% Source revision: 73f18ef013d5b74eafc883efb2c15210642df12a.
% New contribution: complete p=8 nonvanishing, leading Hankel factorization.
\section{The complete order-eight reconstruction theorem}
\label{r8c:section}

The incoming reconstruction report \cite[Section 5]{Reconstruction} proves
optimal-prefix reconstruction for $(p,q)=(8,2b)$ when
$b\not\equiv1\pmod {16}$, while explicitly leaving the remaining
residue class open. We prove nonvanishing for that class as well.
The proof identifies the highest Bernoulli-degree part of the exact
four-dimensional determinant as a positive zeta moment determinant.
A uniform comparison handles all $b\ge512$; exact finite arithmetic
then closes the $31$ outstanding smaller cases. The conclusion is an
exact reconstruction theorem and yields ordinary convergent identities
for products of higher polygammas. No arithmetic independence of
special values is used or claimed.

\subsection{Source coordinates and the matrix whose rank is at issue}

Put
\[
 A_s(n)=\zeta(s,n+1)=\zeta(s)-H_n^{(s)},\qquad x_n=n+\tfrac12,
 \qquad D_{b,r}=D_{8,2b}(-2r),
\]
where
\[
 D_{8,2b}(s)=\sum_{n\ge0}A_8(n)A_{2b}(n)x_n^{-s}
\]
is initially defined for $\operatorname{Re}s>-2b-5$ and is continued
by its centered asymptotic expansion. Its nonpositive even arguments
are regular. The inherited inverse Bernoulli coordinates are
\begin{align}
 C_j&=\sum_{r=0}^j\binom{2j+1}{2r}2^{-2(j-r)}D_{b,r},
 \label{r8c:C}\\
 D_{b,r}&=\sum_{j=0}^r\frac{\binom{2r}{2j}}{2j+1}
 B_{2r-2j}(1/2)C_j.
 \label{r8c:D}
\end{align}
For $j\ge b$ these reduce to a rational linear combination of
$Z_k=\zeta(2k)$, $0\le k<b$, where $Z_0=-1/2$.
The exact late block $M_b=(M_{ik})_{0\le i,k<b}$ is
\begin{equation}\label{r8c:M}
 M_{ik}=B_{2i+2k-6}
 \left\{\frac{\binom{2i+2}{8-2k}}{2i+2}
 +\frac{\binom{2b+2i-6}{2b-2k}}{2b+2i-6}\right\}.
\end{equation}
An entry with negative Bernoulli index is zero. A binomial summand is
omitted unless its lower index is at least $2$ and at most its upper
index. This exclusion includes lower index zero. Thus
\[
 (C_b,\ldots,C_{2b-1})^{\mathsf T}
       =M_b(Z_0,\ldots,Z_{b-1})^{\mathsf T}.
\]
More generally, use the same entry formula for every $i\ge0$ and
$0\le k<b$.  The finite Bernoulli-binomial moment formula from
\cite{GeneralHankel} is
\[
 C_j=\sum_{k=0}^{b-1}M_{j-b,k}Z_k\qquad(j\ge b).
\]
Thus these coefficients also specify every later coordinate,
not only the square block used for the rank proof.
The determinant of the first $2b$ moments in the formal coordinate
list $C_0,\ldots,C_{b-1},Z_0,\ldots,Z_{b-1}$ is
$\det M_b/(4b-1)!!$. This is a statement about rational coefficient
rank, before evaluating the coordinates as complex numbers.

\begin{theorem}[All unequal pairs with smaller order eight]
\label{r8c:main}
For every integer $b\ge5$, $\det M_b\ne0$. Consequently, for every
even $q>8$, the moments
\[
 D_{8,q}(0),D_{8,q}(-2),\ldots,D_{8,q}(-2(q-1))
\]
have full formal rank $q$, and every later centered even moment is
an explicit rational linear combination of this prefix.
\end{theorem}

\subsection{An explicitly defined four-dimensional determinant}

We retain the source normalization, to make the comparison of claims
and certificates unambiguous. For $0\le i,\ell\le3$ set
\[
 n_i=b-4+i,\quad N_{i\ell}=2(n_i-\ell)+1,\quad
 R_\ell=2b-6+2\ell,
\]
\[
 F_{i\ell}=\frac1{(2\ell+1)!}
       \prod_{r=0}^{2\ell-1}(2n_i-r)(2n_i+3+r).
\]
An empty product is one. For
\[
 W_{i\ell}(z)=(z+2\ell+1)
                 (2n_i+2\ell+z+1)^{\underline{2i}},
\]
define the Newton coefficients
\begin{equation}\label{r8c:d}
 d_{i\ell h}=\frac1{h!}\sum_{j=0}^{h}
 (-1)^{h-j}\binom hj W_{i\ell}(N_{i\ell}-j).
\end{equation}
The underline means a falling factorial. With independent variables
$U_1,\ldots,U_7$, put
\begin{equation}\label{r8c:smallA}
 \mathcal A_{i\ell}=2F_{i\ell}+
 \frac{R_\ell!}
 {(2n_i+1)(2n_i+2)(2\ell+2)!N_{i\ell}!}
 \sum_{\substack{1\le h\le2i+1\\ h\ \mathrm{odd}}}
 d_{i\ell h}N_{i\ell}^{\underline h}
 U_{4-i+\ell+(h-1)/2}.
\end{equation}
Ratios of factorials with fixed index difference mean their finite
rational products in $b$.

Set
\begin{align}
 D_8(b)={}&256b^8-3584b^7+18496b^6-40544b^5+25984b^4
 \notag\\
 &+18424b^3-5856b^2-8316b-2025,
 \label{r8c:D8}\\
 P(b;U)&=576(2b-5)D_8(b)\frac{\det\mathcal A(b;U)}{\det F(b)}.
 \label{r8c:Pdef}
\end{align}
The source's finite polynomial certificate proves
$P\in\mathbb Z[b,U_1,\ldots,U_7]$, of degree at most four in
$U$, and gives
\begin{align}
 \det F={}&\frac{256}{4725}(b-2)(b-1)(2b-5)(2b-3)
 (4b-11)(4b-9)\notag\\
 &\hspace{8mm}\cdot(4b-7)^2(4b-5)(4b-3)D_8(b).
 \label{r8c:Fdet}
\end{align}
The integer $D_8(b)$ is odd at every integer $b$, so these denominators
are nonzero in the stated range.

Here is its exact connection to the source matrix. Let
\[
 \mu_r=\begin{cases}B_{2r}/(2r)!,&r\ge0,\\0,&r<0,\end{cases}
 \qquad H_b=(\mu_{i+j-3})_{0\le i,j<b},
\]
\[
 R_b=\operatorname{diag}_{0\le i<b}(2b+2i-7)!,\qquad
 E_b=\operatorname{diag}_{0\le k<b}\frac1{(2b-2k)!}.
\]
For every integer $b\ge8$ the inherited reduction is
\begin{equation}\label{r8c:ratio}
 \frac{\det M_b}{\det R_b\det E_b\det H_b}
 =\frac{P(b;2B_{2b-2},\ldots,2B_{2b-14})}
 {576(2b-5)D_8(b)}.
\end{equation}
The matrix $H_b$ is nonsingular. One short proof conjugates it by
$\operatorname{diag}(1,-1,1,-1,\ldots)$ and partitions after the
first two indices. The lower-right block is a positive definite
moment matrix on the points $(2\pi m)^{-2}$, $m\ge1$. The
upper-left block is zero, and the rectangular off-diagonal block
has row rank two: its first two columns form a nonsingular
antitriangular matrix. Its Schur complement is negative definite.
Hence $H_b$ has exactly two negative eigenvalues and is nonsingular.
Thus nonvanishing of the numerator of \eqref{r8c:ratio} proves the
rank assertion.

All formulas in this subsection are inherited definitions or
reductions. The proof below supplies the previously missing global
nonvanishing. In particular, the symbolic certificate is used as a
finite algebraic identity, not as an extrapolation from values of $b$.

\subsection{The highest-degree part is a zeta moment determinant}

Write $P=P_4+P_{\le3}$ according to total degree in $U_1,\ldots,U_7$.
For an integer $b\ge8$, set
\[
 F_0=(2b-14)!,\qquad
 \alpha_j=\frac{(2b-2j)!}{F_0},\qquad V_j=\frac{U_j}{\alpha_j}
       \quad(1\le j\le7),
\]
and define
\begin{align}
 G(b)={}&12288(b-6)^4(b-5)^4(b-4)^4(b-3)^3(b-2)^2(b-1)
 \notag\\
 &\cdot(2b-13)^4(2b-11)^4(2b-9)^4(2b-7)^3
 (2b-5)^3(2b-3)^2(2b-1).
 \label{r8c:G}
\end{align}

\begin{lemma}[Highest-degree Hankel factorization]\label{r8c:hankel}
The exact identity is
\begin{equation}\label{r8c:P4}
 P_4(b;U)=G(b)\det(V_{i+j+1})_{0\le i,j\le3}.
\end{equation}
At the actual doubled Bernoulli values, this becomes
\begin{equation}\label{r8c:P4zeta}
 P_4(b;2B_{2b-2},\ldots,2B_{2b-14})
 =\frac{256G(b)F_0^4}{(2\pi)^{8b-32}}
       \det\bigl(\zeta(2b-14+2i+2j)\bigr)_{0\le i,j\le3}>0.
\end{equation}
\end{lemma}

\begin{proof}
This factorization has a direct structural proof. In
\eqref{r8c:d}, substitute $z=N_{i\ell}-y$. Then
\[
 W_{i\ell}(N_{i\ell}-y)
    =(2n_i+2-y)(4n_i+2-y)^{\underline{2i}},
\]
which is independent of $\ell$. Write its Newton coefficients as
$d_{ih}$. In particular,
\begin{equation}\label{r8c:d1}
 d_{i1}=-(i+1)(4n_i+2)^{\underline{2i}}.
\end{equation}
Indeed the coefficient is the value at $y=1$ minus the value at
$y=0$; using
$(x-1)^{\underline{2i}}=x^{\underline{2i}}(x-2i)/x$
and $4n_i+2=2(2n_i+1)$ gives \eqref{r8c:d1}.

Let $T=\mathcal A-2F$. In a term with $h=2k+1$, the Bernoulli
variable index is $j=4-i+\ell+k$ and
$N_{i\ell}-h=2b-2j$. Thus substitution of $U_j=\alpha_jV_j$
cancels the factorials in \eqref{r8c:smallA}. Define
\[
 \rho_i=\frac1{(2n_i+1)(2n_i+2)},\qquad
 \sigma_\ell=\frac{R_\ell!}{F_0(2\ell+2)!},
\]
and the lower triangular matrix $L$ by
$L_{i,i-k}=d_{i,2k+1}$ for $0\le k\le i$.
Then the exact matrix identity is
\[
 T=\operatorname{diag}(\rho_i)\,L\,
       (V_{4-i+\ell})_{0\le i,\ell\le3}
       \operatorname{diag}(\sigma_\ell).
\]
Reversing its four middle rows has sign $(-1)^6=1$ and gives
the Hankel matrix in \eqref{r8c:P4}. The diagonal entries of $L$
are \eqref{r8c:d1}. Taking determinants, multiplying by
$576(2b-5)D_8(b)/\det F(b)$, and cancelling the finite products
in \eqref{r8c:Fdet} gives exactly \eqref{r8c:G}.
Since the highest-degree part of $\det(2F+T)$ is $\det T$,
this proves \eqref{r8c:P4}. An independent coefficient check in the
accompanying script verifies all sixteen Hankel monomials.

Euler's Bernoulli formula \cite[Eq. 24.8.1]{DLMF} gives
\[
 V_j=(-1)^{b-j+1}\frac{4F_0}{(2\pi)^{2b-2j}}
                     \zeta(2b-2j).
\]
Factoring the row and column signs and powers of $2\pi$ in the
determinant, and then reversing both its rows and its columns,
gives \eqref{r8c:P4zeta}. All signs cancel. Its positivity follows
from the strictly positive moment representation proved next.
\end{proof}

\begin{lemma}[An explicit positive contribution]\label{r8c:positive}
For every real $n>1/2$,
\begin{equation}\label{r8c:zeta-lower}
 \det\bigl(\zeta(2n+2i+2j)\bigr)_{0\le i,j\le3}
 \ge \frac{30625}{48922361856}\,24^{-2n}>0.
\end{equation}
\end{lemma}
\begin{proof}
With $x_m=m^{-2}$, the entries are the moments
$\sum_{m\ge1}x_m^{n+i+j}$. Cauchy--Binet for a finite truncation,
followed by convergence of each entry, yields
\[
 \det\bigl(\zeta(2n+2i+2j)\bigr)
 =\sum_{m_1<m_2<m_3<m_4}
   (m_1m_2m_3m_4)^{-2n}
   \prod_{i<j}(m_i^{-2}-m_j^{-2})^2.
\]
Every summand is positive. Keeping the term $(1,2,3,4)$ proves
\eqref{r8c:zeta-lower}, since
\[
 \prod_{1\le i<j\le4}(i^{-2}-j^{-2})^2
       =\frac{30625}{48922361856}=:V.
\]
No limiting determinant sign is inferred numerically.
\end{proof}

\subsection{Uniform domination of every lower-degree term}

The integer polynomial defined by \eqref{r8c:Pdef} has $94$
monomials when grouped by the seven $U$ variables. Every coefficient
of $P_{\le3}$ has degree at most $15$ in $b$, and its full integer
coefficient norm is exactly
\begin{equation}\label{r8c:K}
 K:=\sum_{\substack{\text{monomials }c b^rU^\nu\text{ of }P\\
                         |\nu|\le3}}|c|
       =1465775272343372.
\end{equation}
This finite integer is obtained by exact coefficient extraction,
independently replayed from the source polynomial certificate.
It is an upper bound on the whole lower-degree contribution, with
no cancellation assumed.

For $b\ge512$, let
\[
 T_b=\frac{8(2b)^{12}(2b-14)!}{(2\pi)^{2b-14}}.
\]
Every actual variable satisfies $|2B_{2b-2j}|\le T_b$:
use $\zeta(2b-2j)<2$, bound the factorial ratio by $(2b)^{12}$,
and use $2\pi>1$ for the remaining nonnegative exponent.
Also $T_b\ge1$. For example $20!>8^{20}$, the ratio for subsequent
even factorials is greater than $8^2$, and $2\pi<8$.
Consequently
\begin{equation}\label{r8c:small-bound}
 |P_{\le3}(b;2B_{2b-2},\ldots,2B_{2b-14})|
 \le Kb^{15}T_b^3.
\end{equation}
For $b\ge14$, every $b-j$ factor in \eqref{r8c:G} is at least
$b/2$, and every $2b-j$ factor is at least $b$. There are $18$
factors of the first kind and $21$ of the second, with multiplicity.
Thus
\[
 G(b)\ge\frac{3}{64}b^{39}.
\]
Combining this with \eqref{r8c:P4zeta} and
\eqref{r8c:zeta-lower} gives
\[
 P_4\ge
 \frac{12Vb^{39}((2b-14)!)^4}{(2\pi)^{8b-32}}\,24^{-2b+14}.
\]
Dividing \eqref{r8c:small-bound} by this expression, first cancelling
the common powers of $2\pi$ and then using $2\pi<8$, proves
\begin{equation}\label{r8c:ratio-bound}
 \frac{|P_{\le3}|}{P_4}
 \le R_b:=C\,\frac{b^{12}192^{2b}}{(2b-14)!},\qquad
 C=\frac{K2^{73}}{3V24^{14}}
   =\frac{211240479050703917773579966480384}{602791875}.
\end{equation}
The cancellation before replacing $2\pi$ is useful: it avoids
artificially weakening the exponential factor.

All that remains here is elementary exact integer arithmetic:
\begin{equation}\label{r8c:base-bound}
 R_{512}<2^{-675}<1,
 \qquad
 \frac{R_{b+1}}{R_b}
 =\left(1+\frac1b\right)^{12}
   \frac{192^2}{(2b-12)(2b-13)}<1\quad(b\ge512).
\end{equation}
The first inequality means the exact integer comparison
\[
 2^{675}\,211240479050703917773579966480384\,
       512^{12}192^{1024}
       <602791875\cdot1010!,
\]
which is verified in the accompanying script. For the second,
$(513/512)^{12}<2$ and $2\cdot192^2<1011\cdot1012$ suffice.
It follows that $P=P_4+P_{\le3}>0$ for every $b\ge512$.
The threshold is chosen for a transparent uniform proof; no minimal
threshold is asserted.

\subsection{Finite completion of the omitted residue class}

For completeness, the inherited congruence proof gives, for $b\ge11$,
\[
 P(b;2B_{2b-2},\ldots,2B_{2b-14})
       \equiv4(b-1)(4b+15)\pmod {64}
       \quad\text{in }\mathbb Z_{(2)}.
\]
It uses $2B_{2m}\equiv1+4m+16m^2\pmod {64}$ for $m\ge4$,
together with its symbolic determinant certificate. Thus $P\ne0$
when $b\not\equiv1\pmod {16}$. The source checks $5\le b\le10$
directly. Its congruence assertion does not supply the omitted
residue class, and no extrapolation of its valuation formula is
used here.

Below $512$ the only remaining indices are $b=16k+1$,
$1\le k\le31$. The following table records the exact valuation
of $P$ at every such index.
\begin{center}
\begin{tabular}{r r@{\qquad}r r@{\qquad}r r@{\qquad}r r}
\hline
$b$&$v_2(P)$&$b$&$v_2(P)$&$b$&$v_2(P)$&$b$&$v_2(P)$\\
\hline
17&6&33&7&49&6&65&10\\
81&6&97&7&113&6&129&8\\
145&6&161&7&177&6&193&9\\
209&6&225&7&241&6&257&8\\
273&6&289&7&305&6&321&11\\
337&6&353&7&369&6&385&8\\
401&6&417&7&433&6&449&9\\
465&6&481&7&497&6&&\\
\hline
\end{tabular}
\end{center}
Every entry is certified by a nonzero residue modulo $2^{32}$;
all $31$ residues are included in the machine-readable record.
They can be obtained either from exact rational Bernoulli numbers
or entirely in modular arithmetic using
\[
 (2m+1)\,(2B_{2m})
 =2m-1-\sum_{j=1}^{m-1}\binom{2m+1}{2j}\,(2B_{2j}).
\]
The divisor $2m+1$ is odd, so the recurrence is valid modulo
$2^{32}$. Two independent implementations agree on every residue.
In particular no finite numerical approximation is treated as a
nonvanishing certificate. As an explicit first new value,
\[
 \frac{\det M_{17}}{\det R_{17}\det E_{17}\det H_{17}}
   =-\frac{5564518155020169678178976040991}{8937041637533625}.
\]
This settles the previously omitted pair $(p,q)=(8,34)$.
The inherited cases, this finite table, and
\eqref{r8c:base-bound} cover every $b\ge5$, proving
Theorem~\ref{r8c:main}.

\subsection{An all-order family of ordinary polygamma identities}

The rank theorem supplies explicit identities as well as a decision
about a matrix. For $b\ge5$ and $0\le k<b$, define
\begin{align}
 w_{k,r}^{(b)}&=
 \sum_{t=0}^{b-1}(M_b^{-1})_{kt}
 \binom{2(b+t)+1}{2r}2^{-2(b+t-r)},
 \label{r8c:w}\\
 Q_{b,k}(x)&=\sum_{r=0}^{2b-1}w_{k,r}^{(b)}x^{2r},
 \label{r8c:Q}
\end{align}
where the summand is zero for $r>b+t$. These are finite rational
formulas for every coefficient. Put
\[
 \eta_{s,0}=\frac1{s-1},\qquad
 \eta_{s,j}=\frac{(s)_{2j-1}B_{2j}(1/2)}{(2j)!}\quad(j\ge1),
 \qquad
 c_{b,\ell}=\sum_{j=0}^{\ell}\eta_{8,j}\eta_{2b,\ell-j}.
\]
Let $S_{b,k}(x)$ be the polynomial formed by retaining the
nonnegative powers in the formal product
\[
 Q_{b,k}(x)x^{-2b-6}\sum_{\ell\ge0}c_{b,\ell}x^{-2\ell}.
\]
Only finitely many terms enter this polynomial, and its degree is
at most $2b-8$. Then
\begin{equation}\label{r8c:ordinary}
 \boxed{\sum_{n=0}^{\infty}
 \left\{Q_{b,k}(x_n)
 \frac{\psi^{(7)}(n+1)\psi^{(2b-1)}(n+1)}{7!\,(2b-1)!}
       -S_{b,k}(x_n)\right\}=\zeta(2k).}
\end{equation}
This is an ordinary absolutely convergent series. Indeed
\eqref{r8c:w} and \eqref{r8c:C} identify its continued moment
with the $k$th row of $M_b^{-1}M_b$, giving $\zeta(2k)$.
The centered Hurwitz expansion is
\[
 A_8(n)A_{2b}(n)\sim x_n^{-2b-6}
                      \sum_{\ell\ge0}c_{b,\ell}x_n^{-2\ell}.
\]
The polynomial $S_{b,k}$ removes every nondecaying term, leaving
$O(x_n^{-2})$. Each removed nonnegative even power has centered
continued sum $\zeta(-2j,1/2)=0$, so the ordinary convergent sum
has precisely the same value. Thus no unspecified summation of a
divergent series remains in \eqref{r8c:ordinary}.

To reconstruct a later moment, first use \eqref{r8c:C} on the
initial prefix, invert $M_b$, form every later $C_j$ from the
source Bernoulli-binomial formula, and finally use \eqref{r8c:D}.
Every operation is finite rational arithmetic. This completes the
order-eight family. The still-uncovered pairs have smaller even order
at least ten and order gap at least six; their optimal-prefix
conjecture remains separate.

% END EMBEDDED sections/02_rank_eight.tex


% BEGIN EMBEDDED sections/03_eventual_rank.tex
% Optional companion theorem. All labels use ert:.
\section{Effective eventual reconstruction for every fixed even order}
\label{ert:section}

The structural factorization used for order eight has an all-order
consequence. It reduces the full conjecture at each fixed smaller
order to an explicitly bounded finite calculation. It does not prove
that all those finitely many matrices are nonsingular.

\begin{theorem}[Effective eventual optimal-prefix rank]
\label{ert:main}
For every integer $a\ge1$ there is a constructible integer $B_a$
such that the optimal-prefix coefficient matrix $M_{a,b}$ is
nonsingular for every integer $b\ge B_a$. Thus, for every fixed
even tail order $p=2a$, only finitely many larger even orders can
possibly fail the source optimal-prefix conjecture. The formulas
below give a terminating algorithm for an integer $B_a$, using
finite polynomial arithmetic and exact integer comparisons.
\end{theorem}

Here and below, the general source matrix is
\begin{equation}\label{ert:M}
 (M_{a,b})_{ik}=B_{2(i+k-a+1)}
 \left\{\frac{\binom{2i+2}{2a-2k}}{2i+2}
 +\frac{\binom{2b+2i-2a+2}{2b-2k}}{2b+2i-2a+2}\right\},
 \quad0\le i,k<b.
\end{equation}
Negative Bernoulli indices give zero; each binomial summand is
omitted unless its lower index lies between $2$ and its upper index.
This is the same normalization as \eqref{r8c:M} at $a=4$.

\subsection{Exact normalization of the retained determinant}

Assume initially $b\ge2a$. For $0\le i,\ell<a$ define
\[
 n_i=b-a+i,\quad N_{i\ell}=2(n_i-\ell)+1,\quad
 R_\ell=2b-2a+2+2\ell,
\]
\[
 F_{i\ell}=\frac1{(2\ell+1)!}
 \prod_{r=0}^{2\ell-1}(2n_i-r)(2n_i+3+r),
\]
and define integer-polynomial coefficients $d_{ih}$ by
\begin{equation}\label{ert:newton}
 (2n_i+2-y)(4n_i+2-y)^{\underline{2i}}
       =\sum_{h=0}^{2i+1}d_{ih}y^{\underline h}.
\end{equation}
They lie in $\mathbb Z[b]$: the change of basis from ordinary
powers to falling powers has integer Stirling-number coefficients.
Let $U_1,\ldots,U_{2a-1}$ be independent variables and put
\begin{equation}\label{ert:A}
 \mathcal A_{i\ell}=2F_{i\ell}+
 \frac{R_\ell!}
 {(2n_i+1)(2n_i+2)(2\ell+2)!N_{i\ell}!}
 \sum_{k=0}^{i}d_{i,2k+1}N_{i\ell}^{\underline{2k+1}}
            U_{a-i+\ell+k}.
\end{equation}
The inherited reduction is the theorem ``Reduction to a determinant
of size $a$,'' equation \texttt{eq:general-small-det}, in the source
file \texttt{03b\_general\_hankel.tex}~\cite{GeneralHankel}.
We give its exact multiplier and the normalization argument, so
that no nonvanishing statement about an uncomputed inverse matrix
is needed.

\begin{lemma}[Exact source-to-retained determinant multiplier]
\label{ert:normalization}
Put
\[
 h_n=\frac{((2n+1)!)^2}{(4n+2)!(4n+3)!},\qquad
 c_n=\frac{(2n+1)(2n+2)!}{(4n+2)!},
 \qquad \epsilon_{a,b}=a(b-a)+\frac{a(a-1)}2.
\]
For $b\ge2a$, with $U_j=2B_{2b-2j}$,
\begin{equation}\label{ert:exact-multiplier}
 \det M_{a,b}=(-1)^{\epsilon_{a,b}}
 \frac{\prod_{t=0}^{b-1}(2b+2t-2a+1)!}
      {\prod_{k=0}^{b-1}(2b-2k)!}
 \left(\prod_{n=0}^{b-a-1}h_n\right)
 \left(\prod_{i=0}^{a-1}c_{n_i}\right)\det\mathcal A.
\end{equation}
Every scalar factor other than the displayed sign is strictly
positive. In particular,
\begin{equation}\label{ert:equivalence}
 \det M_{a,b}\ne0
 \quad\Longleftrightarrow\quad
 \det\mathcal A\bigl(b;2B_{2b-2},\ldots,2B_{2b-4a+2}\bigr)\ne0.
\end{equation}
\end{lemma}

\begin{proof}
Set $\mu_j=B_{2j}/(2j)!$ for $j\ge0$ and $\mu_j=0$ for $j<0$.
Cancelling factorials in \eqref{ert:M} gives
\[
 M_{a,b}=R(H+U^{\mathrm{pert}})E,
 \quad R_{tt}=(2b+2t-2a+1)!,\quad E_{kk}=\frac1{(2b-2k)!},
\]
where $H_{tk}=\mu_{t+k-a+1}$ and
\[
 U^{\mathrm{pert}}_{tk}=\mathbf1_{k<a}\mu_{t+k-a+1}
 \frac{(2t+1)!(2b-2k)!}
 {(2a-2k)!(2b+2t-2a+1)!}.
\]
The diagonal symbols $R,E$ in this proof are matrices; $R_\ell$
in \eqref{ert:A} is the previously defined integer.

The explicit monic Pad\'e rows are
\[
 \pi_n(x)=\sum_{t=0}^np_{nt}x^t,\qquad
 p_{nt}=\frac{2(2n+2t+1)!(2n+1)!}
 {(4n+2)!(2n-2t+1)!(2t)!}.
\]
For $\mathcal N(x^r)=\mu_{r+1}$ they satisfy
\begin{equation}\label{ert:pade-orthogonality}
 \mathcal N(x^j\pi_n)=0\ (0\le j<n),\qquad
 \mathcal N(x^n\pi_n)=h_n.
\end{equation}
Here is a finite verification of the normalization. Define
\[
 P_N(t)=\sum_{k=0}^N
 \frac{(2N-k)!N!}{(2N)!k!(N-k)!}t^k.
\]
Direct coefficient comparison gives
\[
 e^tP_N(-t)-P_N(t)
 =\frac{(-1)^N(N!)^2}{(2N)!(2N+1)!}t^{2N+1}
       +O(t^{2N+2}).
\]
Indeed, the coefficients through $2N$ reduce to the finite identity
$\sum_k(-1)^k\binom Nk(1+x)^{2N-k}=x^N(1+x)^N$;
the next coefficient is
$(2N)!^{-1}\int_0^1x^N(x-1)^N\,dx$.
Taking $N=2n+1$, set
\[
 q_n(t^2)=\frac{P_N(t)-P_N(-t)}t,\qquad
 r_n(t^2)=\frac{P_N(t)+P_N(-t)}2.
\]
For $\Phi(z)=\sum_{j\ge0}\mu_jz^j
=\frac{\sqrt z}{2}\coth\frac{\sqrt z}{2}$ this proves
\[
 \Phi(z)q_n(z)-r_n(z)=h_nz^{2n+1}+O(z^{2n+2}),
 \qquad \pi_n(x)=x^nq_n(1/x).
\]
Coefficient comparison proves \eqref{ert:pade-orthogonality};
it also gives, for $n\ge\ell\ge0$,
\begin{equation}\label{ert:constant-pade-row}
 \sum_{t=0}^np_{nt}\mu_{t-\ell}=c_n
 \frac1{(2\ell+1)!}
 \prod_{r=0}^{2\ell-1}(2n-r)(2n+3+r).
\end{equation}
Specifically, this sum is $[z^{n-\ell}]r_n(z)$, and its
factorial coefficient formula is the right side. Thus both
$c_n$ and $h_n$ have been obtained without an assumed Pad\'e
normality assertion.

Replace row $n$ of $H+U^{\mathrm{pert}}$ by the linear combination
with coefficients $p_{nt}$, for $0\le n<b$. This simultaneous
transformation is unit lower triangular. In columns $k\ge a$ the
result is $\mathcal N(x^{k-a}\pi_n)$. Hence the final $a$ rows
vanish there, while the first $b-a$ rows in those columns form an
upper triangular square matrix with diagonal
$h_0,\ldots,h_{b-a-1}$. The retained $a$-square matrix is
\[
 A^{\mathrm{ret}}_{ik}=
 \sum_{t=0}^{n_i}p_{n_it}\mu_{t+k-a+1}
 \left(1+\frac{(2t+1)!(2b-2k)!}
 {(2a-2k)!(2b+2t-2a+1)!}\right).
\]
Moving the first $a$ columns past the others contributes
$(-1)^{a(b-a)}$. Reverse the retained columns, putting
$\ell=a-1-k$, and divide row $i$ by $c_{n_i}$.
This adds the determinant factor
$(-1)^{a(a-1)/2}\prod_i c_{n_i}$.

It remains to identify this divided matrix with \eqref{ert:A}.
Its unperturbed entry is $F_{i\ell}$ by
\eqref{ert:constant-pade-row}. Writing $n=n_i$, $N=N_{i\ell}$,
and
\[
 W(j)=(j+2\ell+1)(2n+2\ell+j+1)^{\underline{2i}},\qquad
 C=\frac{R_\ell!}{(2n+1)(2n+2)(2\ell+2)!N!},
\]
factorial cancellation gives the perturbation entry
$2C\sum_{r=0}^{n-\ell}\binom N{2r}B_{2r}W(2r)$.
Since $W(N-y)$ is exactly the polynomial in \eqref{ert:newton},
use the finite identity
\[
 \sum_{j=0}^N\binom Nj B_j(N-j)^{\underline h}
    =N^{\underline h}B_{N-h}(1).
\]
It follows immediately by taking $N^{\underline h}$ outside
the binomial sum and using the defining translation formula for
Bernoulli polynomials. Here $N\ge2i+3$, so all indices $N-h$
are at least two. Only odd $h$ survive. Restoring the omitted
$j=1$ term, with $B_1=-1/2$, gives
\[
 \sum_r\binom N{2r}B_{2r}W(2r)
  =\sum_{k=0}^i d_{i,2k+1}N^{\underline{2k+1}}
           B_{2b-2(a-i+\ell+k)}+\frac N2W(1).
\]
Finally, $CNW(1)=F_{i\ell}$ by factorial cancellation. The divided
retained entry is therefore $2F_{i\ell}$ plus the Bernoulli term
in \eqref{ert:A}. Restoring all row and column factors proves
\eqref{ert:exact-multiplier}.
\end{proof}

For reference, the unperturbed base determinant has the exact
normalization
\begin{equation}\label{ert:base-determinant}
 \det H=(-1)^{\epsilon_{a,b}}
 \left(\prod_{n=0}^{b-a-1}h_n\right)
 \left(\prod_{i=0}^{a-1}c_{n_i}\right)\det F\ne0.
\end{equation}
One can verify the last inequality without using it in the lemma.
Conjugate $H$ by $\operatorname{diag}(1,-1,\ldots)$.
Euler's Bernoulli formula shows that entries with $t+k\ge a$
come from a strictly definite positive moment matrix,
with sign $(-1)^a$. For $a=2h$ the resulting block matrix is
$\left(\begin{smallmatrix}0&B\\B^{\mathsf T}&P\end{smallmatrix}\right)$,
where $P$ is positive definite and $B$ has full row rank $h$.
For $a=2h+1$ it is
$\left(\begin{smallmatrix}ee^{\mathsf T}&B\\B^{\mathsf T}&-P\end{smallmatrix}\right)$,
where $P$ is positive definite, $e=(0,\ldots,0,1)^{\mathsf T}$,
and the first $h$ rows of $B$ have a nonsingular antitriangular
minor. The Schur complements are respectively strictly negative
and strictly positive. The inertia is $(b-h,h)$ in the even case
and $(h+1,b-h-1)$ in the odd case. Together with
\eqref{ert:base-determinant}, this also proves $\det F>0$.
Strict positivity of the moment blocks follows from
$\sum_{m\ge1}(2\pi m)^{-2s}P_0((2\pi m)^{-2})^2>0$ for a
nonzero polynomial $P_0$ and $s\ge1$.

\subsection{The general highest-degree determinant identity}

All denominators can be cleared constructively. Set
\begin{align*}
 J_a(b)&=\prod_{i=0}^{a-1}(2n_i+1)(2n_i+2),\\
 E_a(b)&=\prod_{i,\ell=0}^{a-1}
           \prod_{r=1}^{\max\{0,2i-4\ell-1\}}(R_\ell+r),\\
 D_a(b)&=(2a)!J_a(b)E_a(b),\qquad
 e_a=\sum_{i,\ell=0}^{a-1}\max\{0,2i-4\ell-1\}.
\end{align*}
The symbols $E_a(b)$ and $D_a(b)$ here are auxiliary polynomials;
they are unrelated to the moment coordinates.
Because $N_{i\ell}-R_\ell=2i-4\ell-1$, the polynomial $D_a$
clears every negative factorial ratio in \eqref{ert:A}, every
row denominator, and every constant factorial denominator.
Consequently
\[
 \widetilde{\mathcal A}=D_a\mathcal A
       \in\mathbb Z[b,U_1,\ldots,U_{2a-1}]^{a\times a},
 \qquad \mathscr P_a=\det\widetilde{\mathcal A}
       \in\mathbb Z[b,U_1,\ldots,U_{2a-1}].
\]
Also $D_a(b)>0$ for $b\ge2a$.

Put
\[
 F_0=(2b-4a+2)!,\qquad
 V_j=U_j\frac{F_0}{(2b-2j)!},\qquad
 \kappa_a=\frac{a!((2a)!)^a}{\prod_{\ell=0}^{a-1}(2\ell+2)!},
\]
and
\begin{align}
 \Gamma_a(b)={}&\kappa_aJ_a(b)^{a-1}E_a(b)^a
 \prod_{i=0}^{a-1}(4n_i+2)^{\underline{2i}}
 \prod_{\ell=0}^{a-1}\prod_{r=1}^{2a+2\ell}(2b-4a+2+r).
 \label{ert:Gamma}
\end{align}
The subscript $[a]$ denotes the part homogeneous of total degree
$a$ in the $U$ variables. The exact factorization is
\begin{equation}\label{ert:leading}
 (\mathscr P_a)_{[a]}
 =(-1)^{a(a+1)/2}\Gamma_a(b)
           \det(V_{i+j+1})_{0\le i,j<a}.
\end{equation}
To prove it, write $T=\mathcal A-2F$ and substitute the $V$ variables.
For the term $U_{a-i+\ell+k}$,
$N_{i\ell}-(2k+1)=2b-2(a-i+\ell+k)$, so all nonconstant
factorials cancel. Exactly as in the order-eight proof,
\[
 T=\operatorname{diag}\!\left(\frac1{(2n_i+1)(2n_i+2)}\right)
    L\,(V_{a-i+\ell})_{0\le i,\ell<a}
    \operatorname{diag}\!\left(\frac{R_\ell!}{F_0(2\ell+2)!}\right),
\]
where $L_{i,i-k}=d_{i,2k+1}$ and
$d_{i1}=-(i+1)(4n_i+2)^{\underline{2i}}$.
The $a$ diagonal signs of $L$, together with row reversal of the
middle matrix, give $(-1)^{a(a+1)/2}$. The remaining products are
precisely \eqref{ert:Gamma} after multiplication by $D_a^a$.
This proves \eqref{ert:leading} for every $a$, without evaluating
any large determinant experimentally.

At $U_j=2B_{2b-2j}$, Euler's Bernoulli formula gives
\begin{equation}\label{ert:actual-leading}
 \left|(\mathscr P_a)_{[a]}\right|
 =\frac{\Gamma_a(b)(4F_0)^a}{(2\pi)^{2a(b-a)}}
   \det\bigl(\zeta(2b-4a+2+2i+2j)\bigr)_{0\le i,j<a}.
\end{equation}
Its sign is $(-1)^{ab+a(a+1)/2}$. Define the positive rational
number
\[
 V_a=\prod_{1\le r<s\le a}(r^{-2}-s^{-2})^2,
 \qquad V_1=1.
\]
Keeping the first $a$ support points in Cauchy--Binet proves
\begin{equation}\label{ert:moment-lower}
 \det\bigl(\zeta(2b-4a+2+2i+2j)\bigr)
       \ge V_a(a!)^{-2b+4a-2}>0.
\end{equation}
Thus the highest-degree term is always nonzero in the stated range.

\subsection{An exact algorithm for a sufficient threshold}

Let $\|f\|_1$ denote the sum of the absolute values of all integer
coefficients of a polynomial in $b,U_1,\ldots,U_{2a-1}$. Compute
\[
 E_a^*=\max_{0\le i,\ell<a}\|
                  \widetilde{\mathcal A}_{i\ell}\|_1,
 \qquad K_a=a!(E_a^*)^a,
 \qquad g_a=6a^2-4a+ae_a.
\]
These are finite integer calculations on only $a^2$ explicitly given
entries; no expansion of the whole determinant is required.
The coefficient norm is submultiplicative, so the full coefficient
norm of $\mathscr P_a$ is at most $K_a$.

Each scaled entry has degree at most $6a+e_a-4$ in $b$.
For the nonconstant part this follows from
$\deg_b d_{ih}\le2i+1-h$, the degree $h$ of
$N_{i\ell}^{\underline h}$, and the degree difference
$R_\ell-N_{i\ell}=1+4\ell-2i$; the total bound is
$\deg D_a+4\ell$. The same bound holds for $D_aF_{i\ell}$.
Hence $\deg_b\mathscr P_a\le g_a$.
The factored polynomial \eqref{ert:Gamma} also has degree exactly
$g_a$. Every linear factor in it is at least $b$ when $b\ge4a$,
so
\[
 \Gamma_a(b)\ge\kappa_a b^{g_a}\quad(b\ge4a).
\]

For $b\ge4a+10$, define
\[
 T_{a,b}=\frac{8(2b)^{4a-4}F_0}{(2\pi)^{2b-4a+2}}\ge1.
\]
For the last inequality, put $s_0=2b-4a+2\ge26$;
$26!>8^{26}$ implies $s_0!>8^{s_0}$ by induction,
and $2\pi<8$.
As before, $|2B_{2b-2j}|\le T_{a,b}$ for $1\le j\le2a-1$.
Since the remaining terms have total $U$ degree at most $a-1$,
\[
 |(\mathscr P_a)_{[\le a-1]}|
      \le K_a b^{g_a}T_{a,b}^{a-1}.
\]
Combining this with \eqref{ert:actual-leading} and
\eqref{ert:moment-lower}, cancelling the common powers of $2\pi$,
and then using $1<2\pi<8$, gives
\begin{equation}\label{ert:bound}
 \frac{|(\mathscr P_a)_{[\le a-1]}|}{|(\mathscr P_a)_{[a]}|}
 \le R_a(b):=C_a\frac{b^{L_a}(8a!)^{2b}}{(2b-4a+2)!},
\end{equation}
where every constant is explicit and rational:
\begin{align}
 L_a&=4(a-1)^2,\qquad \chi_a=2a^2-6a+2,\notag\\
 C_a&=\frac{K_a\,8^{a-1+\max(\chi_a,0)}2^{L_a}}
 {\kappa_a4^aV_a(a!)^{4a-2}}.
 \label{ert:constant}
\end{align}
The use of $8^{\max(\chi_a,0)}$ also covers $a=1,2$, where
$\chi_a$ is negative.

Here is a terminating purely arithmetic prescription. Start at
\[
 B_a^{(0)}=\max\{4a+10,\ L_a,\ 16a!\},
\]
and increase the integer $b$ until the exact rational comparison
$R_a(b)<1$ holds. Define $B_a$ to be that first integer. The procedure
must terminate, because for every $b\ge B_a^{(0)}$,
\[
 \frac{R_a(b+1)}{R_a(b)}
 =\left(1+\frac1b\right)^{L_a}
 \frac{(8a!)^2}{(2b-4a+4)(2b-4a+3)}<\frac34.
\]
Indeed $b\ge L_a$ gives $(1+1/b)^{L_a}<3$,
$b\ge4a$ makes both denominator factors at least $b$, and
$b\ge16a!$ supplies the last numerical factor.
The comparison uses integer factorials and powers after clearing
the fixed rational denominator of $C_a$.

For every $b\ge B_a$, \eqref{ert:bound} is strictly less than one,
so the highest-degree term cannot be cancelled. Equations
\eqref{ert:equivalence} and $D_a(b)>0$ prove
Theorem~\ref{ert:main}. In particular, for each input $a$ one may
inspect the finite set $a<b<B_a$ exactly: this either proves the
entire fixed-order family or produces a genuine counterexample.
The construction does not perform those finite calculations for all
$a$, and does not settle the unrestricted conjecture.

\begin{table}[htbp]
\centering
\begin{tabular}{rrrr}
\toprule
$a$ & Smaller order $p=2a$ & Sufficient $B_a$ & Guaranteed even $q$ \\
\midrule
1 & 2  & 16   & $q\ge32$ \\
2 & 4  & 45   & $q\ge90$ \\
3 & 6  & 142  & $q\ge284$ \\
4 & 8  & 467  & $q\ge934$ \\
5 & 10 & 1920 & $q\ge3840$ \\
\bottomrule
\end{tabular}
\caption{Exact sufficient thresholds returned by the constructive
algorithm. No minimality is asserted.}
\label{ert:threshold-table}
\end{table}

The first four rows of Table~\ref{ert:threshold-table} are exact
diagnostics of the general construction; complete-family results
already establish reconstruction for every larger even order in those
cases, including the order-eight theorem proved here. The fifth row
gives a concrete further consequence: for smaller order $p=10$,
optimal-prefix reconstruction holds for every even $q\ge3840$.
The algorithm computes $B_5=1920$ by an exact rational comparison.
It does not check the remaining finite range $6\le b<1920$.

% END EMBEDDED sections/03_eventual_rank.tex


% BEGIN EMBEDDED sections/04_quartic.tex
% Article-ready section.  Requires amsmath, amssymb, amsthm, booktabs.
% Bibliography items appear in quartic_bibliography.tex.
% New claim: the all-moment quartic extension, not the classical zeroth sum.

\section{All regular centered quartic trigamma moments}
\label{qtm:section}

The incoming reconstruction report asks for an explicit quartic extension
of its centered cubic moments, with a fixed coordinate space and a
generator \cite[Question 5 and Section 6]{Reconstruction}.  We answer that
question.  The starting, unweighted quartic sum is already evaluated
explicitly by Hoffman \cite[Corollary 4.4 and the following example]{Hoffman}.
The contribution here is the complete family of centered
polynomial-subtracted moments, the finite Bernoulli formula connecting it
to the cubic family, and the generator with all initial constants fixed.

\subsection{Definition and ordinary summation convention}

Throughout this section set
\begin{equation}\label{qtm:definitions}
 A_n=\psi'(n+1)=\zeta(2)-H_n^{(2)},\qquad
 x_n=n+\tfrac12,\qquad
 \mathcal V(s)=\sum_{n=0}^{\infty}A_n^4x_n^{-s}
 \quad(\Re s>-3).
\end{equation}
Use $B_1=-1/2$ for ordinary Bernoulli numbers.  Define
\begin{equation}\label{qtm:h}
 h_k=\sum_{a+b+c+d=k}
 B_{2a}(1/2)B_{2b}(1/2)B_{2c}(1/2)B_{2d}(1/2),
 \qquad a,b,c,d\ge0.
\end{equation}
For example
\[
 h_0=1,\qquad h_1=-\frac13,\qquad h_2=\frac{19}{120},
 \qquad h_3=-\frac{1871}{15120}.
\]

\begin{proposition}[Regular parity and prescribed subtractions]
\label{qtm:regularity}
The only poles of $\mathcal V$ are simple poles at $s=-3-2k$,
$k\ge0$, with residues $h_k$.  In particular all the values
\begin{equation}\label{qtm:Vr}
 V_r=\mathcal V(-2r),\qquad r\ge0,
\end{equation}
are regular.  They are represented by the ordinary absolutely convergent
series
\begin{equation}\label{qtm:ordinary}
 \boxed{\displaystyle
 V_r=\sum_{n=0}^{\infty}\left(
 x_n^{2r}A_n^4-\sum_{k=0}^{r-2}h_kx_n^{2r-4-2k}\right).}
\end{equation}
The inner sum is empty for $r=0,1$.
\end{proposition}

\begin{proof}
The centered Hurwitz asymptotic expansion gives
\[
 A_n\sim\sum_{j\ge0}B_{2j}(1/2)x_n^{-2j-1},\qquad
 A_n^4\sim\sum_{k\ge0}h_kx_n^{-4-2k}.
\]
Subtract an arbitrary finite number of terms and add the corresponding
Hurwitz functions $h_k\zeta(s+4+2k,1/2)$.  The remainder series is
normally convergent in successively larger left half-planes.  This
continues $\mathcal V$ meromorphically and gives the listed poles and
residues.  Every term in the definition of $h_k$ has sign $(-1)^k$,
because $B_{2j}(1/2)$ has sign $(-1)^j$, including $j=0$.
Consequently $h_k\ne0$ and the poles are actual.

At $s=-2r$ the added terms corresponding to $0\le k\le r-2$
are $h_k\zeta(4+2k-2r,1/2)=0$, since
$\zeta(-2m,1/2)=0$ for every integer $m\ge0$.  The remaining
summand in \eqref{qtm:ordinary} is $O(x_n^{-2})$; for $r=0$ it
is even $O(x_n^{-4})$.  This proves the asserted ordinary identity.
\end{proof}

\subsection{The finite formula}

We use two explicitly evaluated starting constants:
\begin{align}
 Q_5&=4\zeta(2)\zeta(3)-10\zeta(5),\label{qtm:Q5}\\
 Q_7&=51\zeta(3)\zeta(4)+10\zeta(2)\zeta(5)
                  -\frac{301}{4}\zeta(7).\label{qtm:Q7}
\end{align}
The cubic coefficients inherited from the incoming report are
\begin{align}
 C_1&=9\zeta(3)-\frac{17}{4}\zeta(4),\label{qtm:C1}\\
 C_j&=9B_{2j-2}\zeta(3)-\frac12\delta_{j,2}\zeta(2)
       +R_j,\qquad j\ge2,\label{qtm:Cj}\\
 R_j&=-\frac{3}{2(2j-1)}
 \sum_{h=0}^{j-2}(2h+1)\binom{2j-1}{2h+3}
 B_{2j-4-2h}B_{2h}+\frac14\delta_{j,3}.
 \label{qtm:Rj}
\end{align}
These formulas are reproduced to make the new quartic answer directly
computable; their cubic evaluation is not a new claim of this section.

Define $D_0=Q_7$ and, for $j\ge1$, set
\begin{equation}\label{qtm:Dj}
 \boxed{\displaystyle
 D_j=\frac2j\sum_{\ell=1}^{j}
 \binom{2j}{2\ell}B_{2j-2\ell}C_\ell+\varepsilon_j,\qquad
 \varepsilon_j=
 \begin{cases}
 Q_5,&j=1,\\
 -3\zeta(3),&j=2,\\
 \dfrac{2j-5}{2}B_{2j-6},&j\ge3.
 \end{cases}}
\end{equation}
In particular
\begin{align}
 D_1&=2C_1+Q_5,\notag\\
 D_2&=\frac{15}{2}\zeta(3)-\frac{17}{4}\zeta(4)
                 -\frac12\zeta(2)-\frac12,\notag\\
 D_3&=-\frac7{10}\zeta(3)+\frac{17}{12}\zeta(4)
                 -\frac56\zeta(2)-\frac35,\notag\\
 D_4&=\frac{23}{35}\zeta(3)-\frac{17}{12}\zeta(4)
                 +\frac7{12}\zeta(2)-\frac{81}{560}.
 \label{qtm:Dexamples}
\end{align}

\begin{theorem}[Complete quartic Bernoulli formula]
\label{qtm:main}
For every $r\ge0$,
\begin{equation}\label{qtm:mainformula}
 \boxed{\displaystyle
 V_r=\frac1{2r+1}\sum_{j=0}^{r}
 \binom{2r+1}{2j+1}B_{2r-2j}(1/2)D_j.}
\end{equation}
Consequently every sum \eqref{qtm:ordinary} lies in the fixed space
\begin{equation}\label{qtm:space}
 \operatorname{span}_{\mathbb Q}
 \{Q_7,Q_5,\zeta(2),\zeta(3),\zeta(4),1\}.
\end{equation}
The first six moments span all six formal coordinates: in the order
displayed in \eqref{qtm:space}, their coefficient determinant is
\begin{equation}\label{qtm:det}
 \det [V_0,\ldots,V_5]_{\mathrm{coordinates}}
       =\frac{7361}{49896000}\ne0.
\end{equation}
This is a statement about the specified coefficient formulas, not an
assertion of arithmetic independence of these six real constants.
\end{theorem}

Here are the first four identities in compact form:
\begin{align}
 V_0&=Q_7,\label{qtm:V0}\\
 V_1&=\frac{Q_5}{3}-\frac{Q_7}{12}
                +6\zeta(3)-\frac{17}{6}\zeta(4),\label{qtm:V1}\\
 V_2&=-\frac{Q_5}{6}+\frac{7Q_7}{240}-\frac1{10}\zeta(2)
        -\frac32\zeta(3)+\frac{17}{30}\zeta(4)-\frac1{10},
        \label{qtm:V2}\\
 V_3&=\frac{7Q_5}{48}-\frac{31Q_7}{1344}+\frac1{168}\zeta(2)
        +\frac{13}{20}\zeta(3)+\frac{17}{672}\zeta(4)
        +\frac{11}{280}.\label{qtm:V3}
\end{align}
The first two are convergent without subtractions.  The next two are,
respectively,
\[
 \sum_{n\ge0}(x_n^4A_n^4-1),\qquad
 \sum_{n\ge0}(x_n^6A_n^4-x_n^2+1/3).
\]

\subsection{Finite summation proof and the boundary cancellation}

Write $\operatorname{CT}_N$ for the coefficient of
$N^0(\log N)^0$ in a finite-cutoff asymptotic expansion.  It will be
applied only after an exact finite identity has been established.
Set
\begin{equation}\label{qtm:T}
 T_{d,m}=\operatorname{CT}_N\sum_{k=1}^Nk^m A_k^d,
 \qquad b_d(a)=[N^{-a}]A_N^d.
\end{equation}
The coefficients $b_d(a)$ are rational and are fixed by
\begin{equation}\label{qtm:uncentered}
 A_N\sim\sum_{a\ge1}B_{a-1}N^{-a}.
\end{equation}
When a sum converges, its constant term is its ordinary value.  In
particular, with the descending-index convention for multiple zeta values,
\begin{equation}\label{qtm:linearnegative}
 T_{1,-p}=\zeta(2,p)\quad(p\ge1),\qquad
 T_{0,-p}=\zeta(p)\quad(p\ge2),\qquad T_{0,-1}=\gamma.
\end{equation}
The following elementary formula will also be used:
\begin{equation}\label{qtm:oddlinear}
 T_{1,q}=\frac{2-q}{4}B_{q-1}-\frac12\delta_{q,1}\gamma,
 \qquad q\ge1\text{ odd}.
\end{equation}
For completeness let $F_q(N)=\sum_{k=1}^Nk^q$.  Finite summation by
parts gives
\[
 \sum_{k=1}^Nk^qA_k
   =F_q(N)A_N+\sum_{k=1}^N\frac{F_q(k-1)}{k^2}.
\]
Insert Faulhaber's polynomial and \eqref{qtm:uncentered}.  For odd
$q\ge3$, only its quadratic and degree-$q$ terms can contribute
nonzero constant coefficients, and they give $(2-q)B_{q-1}/4$.
At $q=1$, the second sum is $(N-H_N)/2$ and the first term has
constant coefficient $1/4$, giving $1/4-\gamma/2$.

The inherited cubic coefficient has the finite-cutoff form
\begin{equation}\label{qtm:Ccutoff}
 C_\ell=3T_{2,2\ell-2}+3T_{1,2\ell-4}
          +T_{0,2\ell-6}+b_3(2\ell).
\end{equation}
This is the cubic summation identity proved in the incoming report.
One may verify its explicit evaluation
\eqref{qtm:C1}--\eqref{qtm:Rj} by applying the preceding summation
by parts once more to the quadratic tail.  Its degree-$2\ell-2$
Faulhaber term cancels the explicit linear tail in
\eqref{qtm:Ccutoff}; the remaining odd linear powers are evaluated
by \eqref{qtm:oddlinear}.  The terms at zero and negative powers
give $C_1$, the $-\zeta(2)/2$ exception at $\ell=2$, and the
boundary contribution $1/4$ at $\ell=3$.  Thus every low-index
exception retained in \eqref{qtm:Cj} is required.

The quartic coordinate corresponding to the odd monomial $k^{2j+1}$ is
\begin{equation}\label{qtm:Dcutoff}
 \widehat D_j=4T_{3,2j-1}+6T_{2,2j-3}
              +4T_{1,2j-5}+T_{0,2j-7}+b_4(2j+1).
\end{equation}
We now prove $\widehat D_j=D_j$ with no limit of a divergent
series taken before the cancellation.

For $j\ge1$, put $m=2j-1$ and
\begin{equation}\label{qtm:Faulhaber}
 F_m(N)=\sum_{k=1}^Nk^m
   =\sum_{\ell=1}^j f_{2\ell}N^{2\ell}+\frac12N^m,
 \qquad
 f_{2\ell}=\frac1{2j}\binom{2j}{2\ell}B_{2j-2\ell}.
\end{equation}
This includes $j=1$.  Since $F_m(k-1)=F_m(k)-k^m$ and
$A_{k-1}=A_k+k^{-2}$, the exact finite identity is
\begin{equation}\label{qtm:cubicSBP}
 \sum_{k=1}^Nk^m A_k^3
 =F_m(N)A_N^3+
 \sum_{k=1}^NF_m(k-1)
       \left(\frac{3A_k^2}{k^2}+\frac{3A_k}{k^4}+\frac1{k^6}\right).
\end{equation}
Multiply this by four and add the other terms in
\eqref{qtm:Dcutoff}.  The coefficient $-k^m/2$ in $F_m(k-1)$
cancels the entire quadratic-tail sum $6T_{2,m-2}$.  The even
Faulhaber powers form precisely four copies of
\eqref{qtm:Ccutoff}; their boundary terms cancel the even part of
$4F_m(N)A_N^3$.  Taking constant coefficients only now yields
\begin{equation}\label{qtm:preboundary}
 \widehat D_j=
 4\sum_{\ell=1}^j f_{2\ell}C_\ell
 +2b_3(2j-1)+b_4(2j+1)
 -2T_{1,2j-5}-T_{0,2j-7}.
\end{equation}

The decisive rational identity is
\begin{equation}\label{qtm:boundary}
 2b_3(2j-1)+b_4(2j+1)=
 \begin{cases}0,&j=1,2,\\B_{2j-6},&j\ge3.\end{cases}
\end{equation}
Here is a direct proof.  In a formal variable $u=N^{-1}$, let
$E_0(u)=\sum_{k\ge0}B_{2k}u^{2k}$.  Then
$A=uE_0(u)-u^2/2$.  The odd parts of the third and fourth powers are
\[
 (A^3)_{\mathrm{odd}}=u^3E_0(u)^3+\frac34u^5E_0(u),\qquad
 (A^4)_{\mathrm{odd}}=-2u^5E_0(u)^3-\frac12u^7E_0(u).
\]
The $E_0^3$ terms cancel in the indicated pair of coefficients,
leaving $[u^{2j-6}]E_0(u)$.  This proves \eqref{qtm:boundary}.

For $j=1$, \eqref{qtm:linearnegative} and the classical evaluation
$\zeta(2,3)=9\zeta(5)/2-2\zeta(2)\zeta(3)$ give the correction
$-2\zeta(2,3)-\zeta(5)=Q_5$.  For $j=2$, the correction is
$-2\zeta(2,1)-\zeta(3)=-3\zeta(3)$.  For $j=3$,
\eqref{qtm:oddlinear} makes the two Euler constants cancel:
\[
 1-2(1/4-\gamma/2)-\gamma=\frac12.
\]
For $j\ge4$, the last ordinary-power sum has zero constant
coefficient, and \eqref{qtm:oddlinear} gives
\[
 B_{2j-6}-\frac{7-2j}{2}B_{2j-6}
       =\frac{2j-5}{2}B_{2j-6}.
\]
Substitution in \eqref{qtm:preboundary} proves
\eqref{qtm:Dj} for every $j\ge1$.

For $j=0$, all sums in \eqref{qtm:Dcutoff} converge and the
boundary vanishes.  Sorting the four tail indices, or expanding the
displayed powers, gives
\begin{align}\label{qtm:D0MZV}
 \widehat D_0={}&24\zeta(2,2,2,1)+12\zeta(2,2,3)
              +12\zeta(2,4,1)+12\zeta(4,2,1)\notag\\
             &+4\zeta(2,5)+6\zeta(4,3)+4\zeta(6,1)+\zeta(7)
              =Q_7.
\end{align}
The last evaluation is Hoffman's explicitly printed quartic example;
his sum starts at $n=1$, so adding $A_0^4=\zeta(2)^4$ removes
his $-175\zeta(8)/24$ term.  It is also reconstructed in the
supplementary exact shuffle/stuffle verifier.

Finally,
\begin{equation}\label{qtm:centeredF}
 \sum_{n=0}^{k-1}x_n^{2r}
  =\frac1{2r+1}\sum_{j=0}^{r}
   \binom{2r+1}{2j+1}B_{2r-2j}(1/2)k^{2j+1}.
\end{equation}
Every subtracted even-power sum in \eqref{qtm:ordinary} has zero
constant coefficient as a polynomial in $N$.  Finite summation by
parts using \eqref{qtm:centeredF} therefore gives
\eqref{qtm:mainformula}.  The fixed space follows from
\eqref{qtm:Cj} and \eqref{qtm:Dj}.  The first six coefficient rows,
provided explicitly in the supplementary data, give
\eqref{qtm:det} by rational elimination.  This completes the proof.

\subsection{An analytic generator with fixed initial constants}

For a holomorphic function $f$ define its anchored repeated primitive by
\begin{equation}\label{qtm:Ik}
 (I^kf)(t)=\frac1{(k-1)!}\int_0^t(t-u)^{k-1}f(u)\,du,
 \qquad k\ge1.
\end{equation}
Paths lie in the disk $|t|<2\pi$; the integral is path-independent
there.  The established square and cubic generators are
\begin{align}
 G_2(t)&=\frac{t}{2\sinh(t/2)}
 \left(3\zeta(3)-\frac1{4t}
       \int_0^t(t-u)u^2\coth(u/2)\,du\right),
       \label{qtm:G2}\\
 G_3(t)&=\frac1{2\sinh(t/2)}I^2
 \left(3\cosh(t/2)G_2(t)-\frac{17}{4}\zeta(4)
                 -\frac14\zeta(2)t^2+\frac{t^4}{96}\right).
       \label{qtm:G3}
\end{align}
Both formulas are inherited \cite{Square,Reconstruction}.
In particular $2\sinh(t/2)G_3(t)=\sum_{j\ge1}C_jt^{2j}/(2j)!$.

\begin{theorem}[Quartic moment generator]
\label{qtm:generator}
The exponential generating series
\begin{equation}\label{qtm:G4def}
 G_4(t)=\sum_{r\ge0}V_r\frac{t^{2r}}{(2r)!}
\end{equation}
converges for $|t|<2\pi$ and equals
\begin{equation}\label{qtm:G4}
 \boxed{\displaystyle
 G_4(t)=\frac{Q_7t+Q_5t^3/6
       +4I^2\!\left(\cosh(t/2)G_3(t)\right)
       -2I^4\!\left(\sinh(t/2)G_2(t)\right)}
       {2\sinh(t/2)}.}
\end{equation}
The apparent singularity at zero is removable, with $G_4(0)=Q_7$.
The quotient is an even holomorphic function on the disk.
\end{theorem}

\begin{proof}
Set
\[
 D(t)=\sum_{j\ge0}D_j\frac{t^{2j+1}}{(2j+1)!},\quad
 C(t)=2\sinh(t/2)G_3(t),\quad E(t)=\frac t2\coth(t/2).
\]
The Bernoulli transform \eqref{qtm:mainformula} is exactly
$G_4(t)=D(t)/(2\sinh(t/2))$.  Summing the convolution in
\eqref{qtm:Dj} gives
\begin{align*}
 D''(t)
   &=4\frac{E(t)C(t)}t+Q_5t-\frac12\zeta(3)t^3
       +\frac12I^4(tE(t))\\
   &=4\cosh(t/2)G_3(t)+Q_5t
       -2I^2\!\left(\sinh(t/2)G_2(t)\right).
\end{align*}
To check the last cancellation directly, \eqref{qtm:G2} implies
\[
 \sinh(t/2)G_2(t)=\frac32\zeta(3)t-\frac14I^2(tE(t)).
\]
Thus the cubic polynomial from $-3\zeta(3)$ at $j=2$ cancels
the first term in this last identity after integration.  Integrating
$D''$ twice with $D(0)=0$, $D'(0)=Q_7$ proves
\eqref{qtm:G4}.  All the Bernoulli series and inherited generators
converge on $|t|<2\pi$; the anchored integrals preserve holomorphy.
The numerator is odd and has first term $Q_7t$, while the denominator
is odd with first term $t$.  This proves the stated parity, value, and
removability.  No assertion that $2\pi$ is the exact Taylor radius is
needed here.
\end{proof}

\subsection{Finite polylogarithmic form and a rational extraction}

Formula \eqref{qtm:G4} is an explicit triangular integral generator.
It also gives a finite polylogarithm expression, with depth at most three,
by the following anchored rule.  Put $z=e^{-t}$ and write
$L_{\boldsymbol a}(t)=\operatorname{Li}_{\boldsymbol a}(z)$, with
the argument on the outer index:
\[
 L_{a_1,\ldots,a_d}(t)
   =\sum_{n_1>\cdots>n_d\ge1}
        \frac{e^{-tn_1}}{n_1^{a_1}\cdots n_d^{a_d}}\qquad(t>0).
\]
The complete inherited cubic expression is reproduced here to make
the finite compiler self-contained
\cite[Section 6, finite depth-two generator]{Reconstruction}:
\begin{align}
 C(t)={}&
 \left(9\zeta(4)+6t\zeta(3)-\frac{t^4}{16}\right)L_2(t)
 +\left(12\zeta(3)-\frac{t^3}{2}\right)L_3(t)
       \notag\\
 &-3t^2L_4(t)-12tL_5(t)-\frac{45}{2}L_6(t)
       -\frac32t^2L_{2,2}(t)
       \notag\\
 &-6t\left(L_{3,2}(t)+L_{2,3}(t)\right)
       -9\left(L_{4,2}(t)+L_{2,4}(t)\right)-12L_{3,3}(t)
       \notag\\
 &+\frac{15}{2}\zeta(6)-6\zeta(3)^2+6t\zeta(5)
       +\frac12t^3\zeta(3)+\frac18t^2\zeta(4)
       -\frac1{48}t^4\zeta(2)-\frac{t^6}{1440}.
 \label{qtm:Cpolylog}
\end{align}
Initially this expression is read for real $0<t<2\pi$; the integral
formula \eqref{qtm:G3} specifies its holomorphic continuation through
zero and throughout the disk.  The derivative identities
\[
 \frac d{dt}L_{a+1,\boldsymbol b}=-L_{a,\boldsymbol b},\qquad
 \frac d{dt}L_{1,\boldsymbol a}
       =-\frac{z}{1-z}L_{\boldsymbol a}
\]
show that multiplication by $z/(1-z)$ raises the needed depth by at
most one.  The incoming expression for
$C(t)=2\sinh(t/2)G_3(t)$ is a finite sum of polynomials times
polylogarithms of depth at most two, plus zeta polynomials, as displayed
in \eqref{qtm:Cpolylog}.  Since
\[
 4\cosh(t/2)G_3(t)=2C(t)+4\frac{z}{1-z}C(t),
\]
the new generator has depth at most three after the two anchored
integrations.  The remaining term uses only classical polylogarithms:
\begin{equation}\label{qtm:Tpolylog}
 T(t):=2\sinh(t/2)G_2(t)
 =2t\zeta(3)+3\zeta(4)-\frac{t^4}{48}
   -\frac{t^2}{2}\operatorname{Li}_2(z)
   -2t\operatorname{Li}_3(z)-3\operatorname{Li}_4(z).
\end{equation}
Consequently the numerator of \eqref{qtm:G4} can alternatively be
compiled from the finite expression
\begin{equation}\label{qtm:finitecompiler}
 Q_7t+\frac{Q_5t^3}{6}+2I^2C
       +4I^2\left(\frac{z}{1-z}C\right)-I^4T.
\end{equation}
This asserts a sufficient finite function class, not its minimal depth.

Here is an exact endpoint prescription for that compilation.  If
$F_j'=-F_{j-1}$, an unanchored $k$-fold primitive of $t^mF_0(t)$ is
\begin{equation}\label{qtm:primitivecompiler}
 (-1)^k\sum_{j=0}^m\binom{k+j-1}{j}
      \frac{m!}{(m-j)!}t^{m-j}F_{k+j}(t).
\end{equation}
Subtract its Taylor polynomial of degree $k-1$ at zero to obtain
the anchored primitive.  For the twofold integration of
$zL_{\boldsymbol a}/(1-z)$, the only possibly divergent endpoint
is the first-index-one value when the polynomial degree is zero.
It must be combined before taking a limit.  The convergent constant
\[
 \lambda_{\boldsymbol a}
  =\lim_{t\downarrow0}
   \left(L_{1,\boldsymbol a}(t)
             -\zeta(\boldsymbol a)\operatorname{Li}_1(e^{-t})\right)
\]
is minus the sum of all the convergent terms in the stuffle of
$(1)$ with $\boldsymbol a$, excluding its leading $(1,\boldsymbol a)$
term.  Indeed, for a nonempty word appearing in $C$ its first index
is at least two, so
\[
 \sum_{n\ge1}\frac{H_{n-1}^{(\boldsymbol a)}
                         -\zeta(\boldsymbol a)}{n}
\]
converges absolutely, and dominated convergence identifies it with
the displayed limit.  The finite harmonic stuffle evaluates this
sum, since
$H_N\bigl(H_N^{(\boldsymbol a)}-\zeta(\boldsymbol a)\bigr)\to0$.
Here $H_N^{(\boldsymbol a)}$ uses the same strict descending-index
convention as the multiple zeta values.  Pure zeta-polynomial terms in
\eqref{qtm:Cpolylog} are included by adjoining the empty word:
$L_{\varnothing}=1$, $\zeta(\varnothing)=1$,
$L_{1,\varnothing}=\operatorname{Li}_1(e^{-t})$, and
$\lambda_{\varnothing}=0$.
In particular, the degree-zero empty-word coefficient is
$15\zeta(6)/2-6\zeta(3)^2$.
With this convention the degree-zero coefficients in $C$ satisfy
$\sum c_{0,\boldsymbol a}\zeta(\boldsymbol a)=C(0)=0$.
Thus all the divergent $\operatorname{Li}_1$ endpoints cancel in
\eqref{qtm:finitecompiler}, leaving exactly the prescribed
$\lambda_{\boldsymbol a}$ combination.  This rule fixes the integration
constants rather than leaving an undetermined polynomial.

The six-coordinate theorem also produces ordinary rational evaluations.
For example set
\begin{align*}
 P(x)={}&5233+697660x^2+9382304x^4+28113792x^6\\
        &+21816576x^8+3886080x^{10},\\
 R(x)={}&3886080x^6+20521216x^4+21456896x^2
                         +\frac{940103168}{315}.
\end{align*}
Then
\begin{equation}\label{qtm:rational}
 \boxed{\displaystyle
 \sum_{n=0}^{\infty}\bigl(P(x_n)\psi'(n+1)^4-R(x_n)\bigr)
       =-\frac{775936}{3}.}
\end{equation}
Its summand is $O(x_n^{-2})$.  This is an exact rational elimination
of the six proved coefficient rows, with the same polynomial
subtractions as \eqref{qtm:ordinary}.

\subsection{Verification and the remaining question}

The supplementary program reconstructs each finite-cutoff coordinate
without using \eqref{qtm:Dj}.  It recursively applies Faulhaber's
polynomial and $A_{k-1}=A_k+k^{-2}$, and expands the convergent
negative-power tail sums into ordinary multiple zeta values.  A separate
shuffle/stuffle and duality computation reduces weights through seven.
It agrees with \eqref{qtm:Dj} for $0\le j\le24$, checks the
boundary cancellation at 24 indices, and reconstructs the first 20
centered moments independently.  It also verifies 24 generator
coefficients, the stated six-coordinate determinant, and the rational
identity \eqref{qtm:rational} by exact arithmetic.

For a separate numerical diagnostic, the defining finite sums in
\eqref{qtm:ordinary} were evaluated at 120 decimal working digits and
supplemented by centered Hurwitz tails.  The two choices
$(N,J)=(120,36)$ and $(240,54)$ use $N$ direct summands and retain
$J$ asymptotic coefficients.  For $0\le r\le8$, the largest
second-cutoff residual against the closed formulas is below
$3.6\times10^{-90}$.  The data record both cutoffs and the residuals.
These are numerical diagnostics, not interval certificates; the finite
summation proof establishes the all-index result.

The quartic case of the stated incoming question is therefore answered.
For fifth and higher powers, the useful next target is an analogue of
\eqref{qtm:boundary} that identifies all surviving lower-tail terms
before constants are taken.  It should determine a fixed coordinate
space and a generator at each power.  The present quartic theorem
does not assert that the same six constants suffice for those powers.


% END EMBEDDED sections/04_quartic.tex


% BEGIN EMBEDDED sections/05_bilateral.tex
% Article-ready fragment. Classical references and the source-report citation
% should be wired into the parent bibliography when integrating this file.
\section{Several quadratic factors beyond the first spectral derivative}
\label{msj:sec:main}

The incoming quadratic-spectrum report \cite{Quadratic} proves additivity of the first
spectral derivative for any even polynomial weight and any finite list
of masses. Its explicit next question asks for the second derivative,
the logarithmic cross term, and the limit when two masses coincide.
The following results answer that question with an ordinary convergent
series and a one-dimensional integral of elementary functions. They also
identify the finite corrections needed at every higher spectral order.
The classical multiplicative-anomaly framework \cite{EVZ} concerns first derivatives
of spectral zeta functions; the statements below concern the additional
higher derivatives. No historical-priority claim is needed for this
continuation of the inspected repository.

Fix $0<a<1$, put $v_n=(n+a)^2$ for $n\in\mathbb Z$, and let
\[
 P(v)=\sum_{d=0}^{D}p_dv^d,\qquad B_j>0\quad(1\le j\le r),\qquad r\ge2.
\]
The parameter $B_j$ is the squared mass. Initially in their respective
half-planes of absolute convergence, define
\begin{align}
 \mathcal E_P(s;a,B)&=\sum_{n\in\mathbb Z}P(v_n)(v_n+B)^{-s},
 \label{msj:eq:E}\\
 \mathcal Z_{P,\boldsymbol B}(s;a)
 &=\sum_{n\in\mathbb Z}P(v_n)
                  \prod_{j=1}^{r}(v_n+B_j)^{-s}.
 \label{msj:eq:Z}
\end{align}
Their continuations are holomorphic at $s=0$. Indeed, a tail expansion
uses integer powers of $v_n^{-1}$, and a possible pole at zero would
require an odd power $(n+a)^{-1}$. The detailed subtraction proof below
also supplies this continuation directly.

Write
\begin{equation}
 \mathfrak D_k(P;a,\boldsymbol B)
 =r^{k-1}\sum_{j=1}^{r}\mathcal E_P^{[k]}(0;a,B_j)
       -\mathcal Z_{P,\boldsymbol B}^{[k]}(0;a),
 \qquad k\ge1,
 \label{msj:eq:defect}
\end{equation}
where square brackets indicate derivatives in the spectral parameter.
The factors $r^{k-1}$ are part of the definition. They result from
comparing the factors at the same spectral scale:
\begin{equation}
 \mathfrak D(s)=\frac1r\sum_j\mathcal E_P(rs;a,B_j)
                           -\mathcal Z_{P,\boldsymbol B}(s;a),
 \qquad \mathfrak D_k=\mathfrak D^{[k]}(0).
 \label{msj:eq:common-scale}
\end{equation}
The already established first-derivative identity is $\mathfrak D_1=0$.

\subsection{An exact second-derivative formula}

For a function with a Laurent expansion at $v=\infty$, let
$\operatorname{Pol}_{\infty}$ retain its nonnegative integer powers of
$v$. Define
\begin{align}
 T_P(v;B,C)
 &=\operatorname{Pol}_{\infty}
       \left\{P(v)\log^2\frac{v+B}{v+C}\right\},
 \label{msj:eq:T2}\\
 S_P(a;B,C)
 &=\sum_{n\in\mathbb Z}
 \left\{P(v_n)\log^2\frac{v_n+B}{v_n+C}-T_P(v_n;B,C)\right\}.
 \label{msj:eq:S2}
\end{align}
Every term in the subtraction is explicit. In fact,
\begin{align}
 q_k(B,C)&=(-1)^k\sum_{m=1}^{k-1}
 \frac{(B^m-C^m)(B^{k-m}-C^{k-m})}{m(k-m)},\qquad k\ge2,
 \label{msj:eq:qcoeff}\\
 T_P(v;B,C)&=\sum_{d=2}^{D}p_d\sum_{k=2}^{d}q_k(B,C)v^{d-k}.
 \label{msj:eq:T2-explicit}
\end{align}
Empty sums are zero. Thus the weights $P(v)=1$ and $P(v)=v$ need no
subtraction. For $P(v)=v^2$ the only subtraction is $(B-C)^2$, while
for $P(v)=v^3$ it is
\[
 (B-C)^2v-(B-C)^2(B+C).
\]

\begin{theorem}[Second spectral derivative for any number of masses]
\label{msj:thm:second}
The series \eqref{msj:eq:S2} is ordinary and absolutely convergent, and
\begin{equation}
 \boxed{\displaystyle
 \mathfrak D_2(P;a,\boldsymbol B)
     =\sum_{1\le i<j\le r}S_P(a;B_i,B_j).}
 \label{msj:eq:second-pairwise}
\end{equation}
Consequently
\begin{equation}
 \mathcal Z_{P,\boldsymbol B}^{[2]}(0;a)
 =r\sum_j\mathcal E_P^{[2]}(0;a,B_j)
                       -\sum_{i<j}S_P(a;B_i,B_j).
 \label{msj:eq:second-evaluation}
\end{equation}
For $\deg P\le1$, the pair correction is the unsubtracted convergent
sum of logarithm squares. In particular, for $P=1$ it is nonzero whenever
at least two masses are distinct. Thus first-derivative additivity does
not extend unchanged to second derivatives.
\end{theorem}

\begin{proof}
At a fixed $v>0$ write $L_j=\log(v+B_j)$. Twice differentiating the
common-scale summand at zero gives
\[
 P(v)\left\{r\sum_jL_j^2-\left(\sum_jL_j\right)^2\right\}
        =P(v)\sum_{i<j}(L_i-L_j)^2.
\]
The logarithm difference has the expansion
\[
 \log\frac{v+B}{v+C}
   =\sum_{m\ge1}\frac{(-1)^{m+1}(B^m-C^m)}{m v^m}
\]
on a sufficiently large tail. Squaring proves
\eqref{msj:eq:qcoeff} and \eqref{msj:eq:T2-explicit}; the subtracted
summand is $O(v^{-1})$, uniformly for masses in a compact positive set.
It therefore has an absolutely convergent bilateral sum.

It remains to retain the exact value of the removed polynomial. Put
\begin{equation}
 \mathcal H_a(s)=\zeta(2s,a)+\zeta(2s,1-a).
 \label{msj:eq:H}
\end{equation}
For every integer $m\ge0$,
\begin{equation}
 \mathcal H_a(-m)=0,
 \label{msj:eq:trivial-zero}
\end{equation}
because $\zeta(-2m,x)=-B_{2m+1}(x)/(2m+1)$ and
$B_{2m+1}(1-a)=-B_{2m+1}(a)$. Hence each removed $v^m$ has zero
bilateral spectral value. The common-scale normal-convergence argument
in Theorem~\ref{msj:thm:alljets} makes this differentiation and restoration
precise and gives zero finite correction at order two. This proves
\eqref{msj:eq:second-pairwise}. With an affine weight, the raw tail is
$O(v^{-1})$ or better. With $P=1$ every pair summand is nonnegative and
a distinct pair has strictly positive summands, giving the final claim.
\end{proof}

For two masses, the mixed logarithmic coordinate is equivalently
\begin{equation}
 \operatorname{Reg}\sum_{n\in\mathbb Z}
  P(v_n)\log(v_n+B)\log(v_n+C)
 =\frac{\mathcal E_P^{[2]}(0;a,B)+\mathcal E_P^{[2]}(0;a,C)
                     -S_P(a;B,C)}2.
 \label{msj:eq:cross-term}
\end{equation}
Here $\operatorname{Reg}$ denotes the analytic continuation determined
by the common quadratic spectral coordinate. Formula
\eqref{msj:eq:S2}, or the ordinary integral below, makes the new
cross-term correction explicit; it does not subtract two unspecified
divergent sums.

\subsection{An elementary integral for the entire correction}

The classical bilateral resolvent, from cotangent partial fractions \cite{DLMF}, is
\begin{equation}
 R_a(t)=\sum_{n\in\mathbb Z}\frac1{v_n+t}
       =\frac{\pi}{\sqrt t}
        \frac{\sinh(2\pi\sqrt t)}
             {\cosh(2\pi\sqrt t)-\cos(2\pi a)},\qquad t>0.
 \label{msj:eq:resolvent}
\end{equation}
It follows by pairing the cotangent partial fractions at
$a+i\sqrt t$ and $a-i\sqrt t$. Set
\[
 f_P(t)=P(-t)R_a(t).
\]

\begin{theorem}[Ordinary logarithmic quadrature]
\label{msj:thm:quadrature}
For $B>C>0$,
\begin{equation}
 \boxed{\displaystyle
 S_P(a;B,C)=2\int_C^B f_P(t)
                     \log\frac{B-t}{t-C}\,\mathrm dt.}
 \label{msj:eq:quadrature}
\end{equation}
The integral is ordinary and absolutely convergent; its only endpoint
singularities are logarithmic. For arbitrary positive $B,C$, use
$S_P(a;B,C)=S_P(a;C,B)$ and $S_P(a;B,B)=0$. A form with a continuous
integrand on the square is
\begin{equation}
 S_P(a;B,C)=\int_C^B\int_C^B
      \frac{f_P(t)-f_P(u)}{u-t}\,\mathrm du\,\mathrm dt,
 \label{msj:eq:double-quadrature}
\end{equation}
where the diagonal value is $-f_P'(t)$.
\end{theorem}

\begin{proof}
Polynomial division gives
\[
 P(v)/(v+t)=Q_t(v)+P(-t)/(v+t)
\]
with a polynomial $Q_t$. Therefore subtraction of the polynomial part
and summation of the rational remainder give
\begin{align*}
 &\sum_{n\in\mathbb Z}(1-\operatorname{Pol}_{\infty})
          \frac{P(v_n)}{(v_n+t)(v_n+u)}\\
 &\hspace{35mm}=\frac{P(-t)R_a(t)-P(-u)R_a(u)}{u-t}.
\end{align*}
The statement at $u=t$ follows by continuity. For $(t,u)$ in a compact
positive square, every subtracted summand is uniformly $O(v_n^{-1})$.
Normal convergence permits integration of this identity over $[C,B]^2$.
Since
\[
 \int_C^B\frac{\mathrm dt}{v+t}=\log\frac{v+B}{v+C},
\]
and finite polynomial extraction commutes with these integrals, the
left side becomes \eqref{msj:eq:S2}. This proves
\eqref{msj:eq:double-quadrature}.

To reduce the square integral, omit temporarily the strip
$|u-t|\le\varepsilon$, split the numerator, and interchange $u,t$ in
its second term. The result tends to
\[
 2\int_C^B f_P(t)\bigl(\log(B-t)-\log(t-C)\bigr)\,\mathrm dt.
\]
The omitted-strip contribution tends to zero because the divided
difference extends continuously to the diagonal. The endpoint
logarithms are integrable. Thus the temporary symmetric exclusion is
only a proof device; the final integral needs no principal value or
finite-part convention.
\end{proof}

For example, at $a=1/2$, $P=1$, $B=4$, $C=1$, the exact correction is
\begin{equation}
 \sum_{n\in\mathbb Z}
 \log^2\frac{(n+1/2)^2+4}{(n+1/2)^2+1}
 =2\pi\int_1^4\frac{\tanh(\pi\sqrt t)}{\sqrt t}
                      \log\frac{4-t}{t-1}\,\mathrm dt.
 \label{msj:eq:example-unweighted}
\end{equation}
At the quartic weight $P(v)=v^2$, the corresponding exact identity is
\begin{equation}
 \sum_{n\in\mathbb Z}\left[
 (n+1/2)^4\log^2\frac{(n+1/2)^2+4}{(n+1/2)^2+1}-9\right]
 =2\pi\int_1^4 t^{3/2}\tanh(\pi\sqrt t)
                       \log\frac{4-t}{t-1}\,\mathrm dt.
 \label{msj:eq:example-weighted}
\end{equation}
Both series and both integrals are ordinary convergent expressions.

\subsection{Every confluent coefficient in finite resolvents}

Let $H_m=\sum_{j=1}^m1/j$ and $H_0=0$, and put
$\rho=\min(a,1-a)$. For $k\ge1$, define the rationally subtracted
resolvent moment
\begin{equation}
 R_{P,k}(a;C)
 =\sum_{n\in\mathbb Z}(1-\operatorname{Pol}_{\infty})
                         \frac{P(v_n)}{(v_n+C)^k}.
 \label{msj:eq:Rpk}
\end{equation}
As above, this is an absolutely convergent series after its explicit
polynomial division.

\begin{theorem}[Complete confluent expansion]
\label{msj:thm:confluent}
For $C>0$ and $|h|<C+\rho^2$,
\begin{align}
 S_P(a;C+h,C)
 &=2\sum_{k=2}^{\infty}\frac{(-1)^kH_{k-1}}{k}
                       h^k R_{P,k}(a;C),
 \label{msj:eq:confluent-resolvent}\\
 &=-2\sum_{k=2}^{\infty}\frac{H_{k-1}}{k!}
                       f_P^{(k-1)}(C)h^k.
 \label{msj:eq:confluent-elementary}
\end{align}
Here $S_P(a;C+h,C)$ denotes its principal-logarithm analytic
continuation from positive masses to this disk.  Indeed every
$v_n+C+h$ has positive real part on the disk, so the logarithms
and the explicitly subtracted series define a single holomorphic
function there.
Each coefficient has the finite reduction
\begin{equation}
 R_{P,k}(a;C)
 =\frac{(-1)^{k-1}}{(k-1)!}f_P^{(k-1)}(C)
 =\sum_{\ell=0}^{\min(D,k-1)}
       \frac{P^{(\ell)}(-C)}{\ell!}
                          \mathcal E_1(k-\ell;a,C).
 \label{msj:eq:Rpk-finite}
\end{equation}
The values $\mathcal E_1(N;a,C)$ at positive integers are obtained by
differentiating the elementary function \eqref{msj:eq:resolvent}, or
by the manuscript's finite polygamma partial fractions.
\end{theorem}

\begin{proof}
The elementary Taylor identity
\[
 \log^2(1+x)=2\sum_{k=2}^{\infty}\frac{(-1)^kH_{k-1}}{k}x^k,
 \qquad |x|<1,
\]
follows by differentiating both sides and using
$\log(1+x)/(1+x)=\sum_{m\ge1}(-1)^{m+1}H_mx^m$.
Apply it with $x=h/(v_n+C)$, extract the polynomial parts before
summation, and use locally uniform convergence on the stated disk.
This proves \eqref{msj:eq:confluent-resolvent}.

Differentiating $f_P(C)$ under the subtracted resolvent sum proves the
first equality in \eqref{msj:eq:Rpk-finite}. For the second, expand
$P(v)$ at $v=-C$. Terms with $\ell\ge k$ are polynomials and disappear
under $1-\operatorname{Pol}_{\infty}$; the other terms give the displayed
finite sum. The derivative formula yields
\eqref{msj:eq:confluent-elementary}. The function $f_P$ is meromorphic
in $C$, with possible poles only at $C=-v_n$. The nearest possible pole
is at distance $C+\rho^2$, which also proves the claimed disk of
convergence. Zeros of $P$ can only remove poles and enlarge this disk.
\end{proof}

The first terms are
\[
 S_P(a;C+h,C)=h^2R_{P,2}(a;C)-h^3R_{P,3}(a;C)
                   +\frac{11}{12}h^4R_{P,4}(a;C)+O(h^5).
\]
In particular, for $P=1$, $a=1/2$, and $C=1$,
\begin{equation}
 S_1(1/2;1+h,1)
 =\frac12\left(\pi\tanh\pi-\pi^2\operatorname{sech}^2\pi\right)h^2
                                                    +O(h^3).
 \label{msj:eq:confluent-example}
\end{equation}
This explicitly evaluates the first nonzero coefficient when two
quadratic factors merge.

\subsection{A finite correction formula at every spectral order}

The vanishing \eqref{msj:eq:trivial-zero} removes polynomial
counterterms at the second derivative. Higher derivatives also produce
polynomials times powers of $\log v$, whose spectral values need not
vanish. The following formula retains every such contribution.

Define
\begin{align}
 A(s,z)&=\frac1r\sum_{j=1}^r(1+B_jz)^{-rs}
                         -\prod_{j=1}^r(1+B_jz)^{-s}
          =\sum_{q\ge2}a_q(s)z^q,
 \label{msj:eq:A}\\
 a_q(s)&=(-1)^q\left\{
  \frac1r\sum_{j=1}^r\frac{(rs)_qB_j^q}{q!}
  -\sum_{q_1+\cdots+q_r=q}\prod_{j=1}^r
                      \frac{(s)_{q_j}B_j^{q_j}}{q_j!}\right\}.
 \label{msj:eq:aq}
\end{align}
The coefficient polynomials $a_q(s)$ are divisible by $s^2$; the
coefficients $a_0,a_1$ vanish identically. Put
\begin{equation}
 T(s,v)=\sum_{d=2}^{D}p_d\sum_{q=2}^{d}a_q(s)v^{d-q}.
 \label{msj:eq:T-spectral}
\end{equation}

\begin{theorem}[All-order ordinary remainder and finite Hurwitz jets]
\label{msj:thm:alljets}
For $s$ in a neighborhood of zero,
\begin{align}
 \mathfrak D(s)
 &=\sum_{n\in\mathbb Z}v_n^{-rs}
             \bigl(P(v_n)A(s,v_n^{-1})-T(s,v_n)\bigr)
 \notag\\
 &\quad+\sum_{d=2}^{D}p_d\sum_{q=2}^{d}
                           a_q(s)\mathcal H_a(rs+q-d).
 \label{msj:eq:regulator}
\end{align}
The first series is normally convergent near zero. For $k\ge2$, define
\begin{align}
 F_k(v)&=(-1)^kP(v)\left\{
     r^{k-1}\sum_j\log^k(v+B_j)
            -\left(\sum_j\log(v+B_j)\right)^k\right\},
 \label{msj:eq:raw-k}\\
 T_k(v)&=\sum_{d=2}^{D}p_d\sum_{q=2}^{d}
       \sum_{j=2}^{k}\binom{k}{j}a_q^{[j]}(0)
                (-r\log v)^{k-j}v^{d-q},
 \label{msj:eq:Tk}\\
 C_k&=\sum_{d=2}^{D}p_d\sum_{q=2}^{d}
       \sum_{j=2}^{k-1}\binom{k}{j}a_q^{[j]}(0)
                 r^{k-j}\mathcal H_a^{[k-j]}(q-d).
 \label{msj:eq:Ck}
\end{align}
Then the exact $k$th spectral defect is
\begin{equation}
 \boxed{\displaystyle
 \mathfrak D_k=\sum_{n\in\mathbb Z}\bigl(F_k(v_n)-T_k(v_n)\bigr)+C_k.}
 \label{msj:eq:alljets}
\end{equation}
The remainder series is ordinary and absolutely convergent. Its
summand is $O(v_n^{-1}(1+|\log v_n|^{k-2}))$. All correction terms are
finite Hurwitz-zeta jets at nonpositive even integers. For affine $P$,
both $T_k$ and $C_k$ vanish at every $k\ge2$.
\end{theorem}

\begin{proof}
At $s=0$, the two terms in \eqref{msj:eq:A} both equal one, and their
first $s$ derivatives both equal $-\sum_j\log(1+B_jz)$. Thus
$a_q(s)=O(s^2)$. Expansion to first degree in $z$ shows that $a_1=0$.
The Pochhammer expansion proves \eqref{msj:eq:aq} exactly.

On a sufficiently large lattice tail, Taylor expansion in $z=v^{-1}$
and the definition of $T$ give
\[
 P(v)A(s,v^{-1})-T(s,v)=O(v^{-1})
\]
uniformly for $s$ in a compact neighborhood of zero. Multiplication by
$v^{-rs}$ leaves a normally convergent series, for example on
$|\operatorname{Re}s|<1/(4r)$. Finitely many smaller $v_n$ cause no
difficulty: the exact functions $A(s,v_n^{-1})$ are used there,
without applying a possibly invalid small-$z$ expansion.

Restoring the removed monomials by
$\sum_n v_n^{d-q-rs}=\mathcal H_a(rs+q-d)$ gives
\eqref{msj:eq:regulator}, first in a common half-plane after a finite
head is separated and then by meromorphic continuation. Alternatively,
the same head-tail construction is a direct definition of this
continuation. None of the restored factors has a pole at zero:
$q-d$ is an integer, whereas the only pole of $\mathcal H_a$ is at
$1/2$. This proves the asserted holomorphy and normal convergence.

Differentiate \eqref{msj:eq:regulator} $k$ times. The first raw
derivative is \eqref{msj:eq:raw-k}, and the derivative of
$v^{-rs}T(s,v)$ is \eqref{msj:eq:Tk}. In the restored finite sum,
the terms with $j=0,1$ vanish because $a_q$ is divisible by $s^2$;
the term $j=k$ vanishes by \eqref{msj:eq:trivial-zero}. The remaining
terms are exactly \eqref{msj:eq:Ck}. Derivatives of the normal tail
give the stated polynomial-logarithmic bound. For $D\le1$, all finite
sums are empty.
\end{proof}

The second-derivative theorem follows immediately: at $k=2$ the
finite sum $C_2$ is empty, the raw derivative is the pairwise square,
and $T_2$ is precisely its polynomial part. At the third derivative
the correction can already be nonzero. For two masses and $P(v)=v^2$,
\begin{align}
 F_3(v)&=-3v^2\log\bigl((v+B)(v+C)\bigr)
                         \log^2\frac{v+B}{v+C},
 \label{msj:eq:third-raw}\\
 T_3(v)&=-6(B-C)^2\log v,
 \label{msj:eq:third-subtraction}\\
 C_3&=-12(B-C)^2\log(2\sin\pi a).
 \label{msj:eq:third-contact}
\end{align}
Indeed $a_2(s)=\tfrac12(B-C)^2s^2$ in this case, while Lerch's
identity and Gamma reflection give
\[
 \mathcal H_a'(0)=2\bigl(\zeta'(0,a)+\zeta'(0,1-a)\bigr)
                 =-2\log(2\sin\pi a).
\]
Formula \eqref{msj:eq:third-contact} is therefore a concrete check on
the sign and the scale factor in the general correction law. Discarding
it would change the third spectral derivative.

\subsection{Polylogarithmic coordinates for the finite corrections}

Define the symmetric polylogarithm
\[
 \mathcal C_s(a)=\frac{\operatorname{Li}_s(e^{2\pi ia})
                          +\operatorname{Li}_s(e^{-2\pi ia})}{2}.
\]
For real $s$ this is the real part of the first polylogarithm; the
symmetric definition retains holomorphic dependence on complex $s$.
The symmetric Hurwitz functional equation gives, for $m\ge0$ and
$u$ near zero,
\begin{equation}
 \mathcal H_a(-m+u)
 =4(-1)^m\frac{\Gamma(2m+1-2u)}{(2\pi)^{2m+1-2u}}
                  \sin(\pi u)\mathcal C_{2m+1-2u}(a).
 \label{msj:eq:H-polylog}
\end{equation}
Consequently $\mathcal H_a^{[h]}(-m)$ is a finite combination of
polylogarithm order derivatives at $2m+1$, of order at most $h-1$,
and Gamma derivatives at $2m+1$. The zero of $\sin(\pi u)$ is
responsible for this reduction in order. For example,
\begin{equation}
 \mathcal H_a'(-m)=
 \frac{2(-1)^m(2m)!}{(2\pi)^{2m}}\mathcal C_{2m+1}(a).
 \label{msj:eq:Hprime-polylog}
\end{equation}
Thus the finite corrections in \eqref{msj:eq:Ck} use polylogarithm
order derivatives only through order $k-3$. This is a finite reduction
of the correction, not a claim that the ordinary remainder in
\eqref{msj:eq:alljets} always reduces to finitely many polylogarithms.

\subsection{A confluent law simultaneously at every spectral order}

\begin{theorem}[All-order mass-coalescence law]
\label{msj:thm:allconfluent}
Let $B>0$ and fixed real numbers $c_1,\ldots,c_r$ be given, and set
$B_j=B+\varepsilon c_j$. For every $k\ge2$,
\begin{align}
 \mathfrak D_k(P;a,\boldsymbol B)
 &=\binom{k}{2}r^{k-2}\varepsilon^2
       \left(r\sum_jc_j^2-\Bigl(\sum_jc_j\Bigr)^2\right)
                    \mathcal E_P^{[k-2]}(2;a,B)
       +O(\varepsilon^3).
 \label{msj:eq:allconfluence}
\end{align}
The right side is well defined even when its defining raw series at
$s=2$ diverges: $\mathcal E_P(s;a,B)$ has no pole at an integer.
For $\deg P\le1$ it is an ordinary spectral derivative of a convergent
resolvent series. If all $c_j$ are equal, the masses remain identical
and every defect vanishes exactly.
\end{theorem}

\begin{proof}
For a fixed $v$, expand the common-scale expression in
\eqref{msj:eq:common-scale} through second order in $\varepsilon$.
The constant and linear coefficients cancel. The quadratic term is
\[
 \frac{s^2\varepsilon^2}{2}
 \left(r\sum_jc_j^2-\Bigl(\sum_jc_j\Bigr)^2\right)(v+B)^{-rs-2}.
\]
The same tail subtraction as above justifies this identity for the
continued families, locally uniformly in $s$ near zero. Hence
\[
 \mathfrak D(s)=\frac{s^2\varepsilon^2}{2}
 \left(r\sum_jc_j^2-\Bigl(\sum_jc_j\Bigr)^2\right)
                 \mathcal E_P(rs+2;a,B)+O(\varepsilon^3).
\]
Taking its $k$th derivative gives \eqref{msj:eq:allconfluence}.
The single-mass tail expansion has possible poles only at half-integers,
which proves the assertion at $s=2$. When all masses are equal,
$\mathcal Z_{P,\boldsymbol B}(s;a)=\mathcal E_P(rs;a,B)$, so
\eqref{msj:eq:common-scale} is identically zero.
\end{proof}

\subsection{A global Hurwitz-jet evaluation of the individual terms}

For completeness, the individual second derivatives in
\eqref{msj:eq:second-evaluation} can be evaluated without leaving an
unevaluated multivariate spectral zeta function. Choose an integer
$N\ge0$ such that $B<(N+\rho)^2$, and define
\[
 \mathcal H_{a,N}(s)=\zeta(2s,N+a)+\zeta(2s,N+1-a).
\]
Separating the $2N$ lattice terms with positive coordinates
$n+a$ and $n+1-a$, $0\le n<N$, gives the geometrically convergent
formula
\begin{align}
 \mathcal E_P^{[2]}(0;a,B)
 &=\sum_{n=0}^{N-1}\sum_{x\in\{n+a,n+1-a\}}
                         P(x^2)\log^2(x^2+B)
 \notag\\
 &\quad+\sum_{d=0}^{D}p_d\Bigg\{
    \mathcal H_{a,N}^{[2]}(-d)
 \notag\\
 &\hspace{20mm}+2\sum_{q=1}^{\infty}\frac{(-1)^qB^q}{q}
       \left[H_{q-1}\mathcal H_{a,N}(q-d)
                         +\mathcal H_{a,N}'(q-d)\right]\Bigg\}.
 \label{msj:eq:individual-second}
\end{align}
Every spectral argument here is an integer in the $\mathcal H$ variable,
so no pole is substituted separately. This identity follows directly
from
\[
 (s)_q=(q-1)!s\bigl(1+H_{q-1}s+O(s^2)\bigr)
\]
and the binomial expansion of $(x^2+B)^{-s}$ on the chosen tail.
The strict tail-radius condition justifies spectral differentiation
term by term. Higher derivatives use the same finite Pochhammer
coefficient calculus, or unsigned Stirling numbers of the first kind.
Together with \eqref{msj:eq:quadrature}, this supplies a global exact
evaluation of the several-mass second derivative in ordinary Hurwitz
jets and an explicitly anchored elementary integral.

\subsection{Verification and the remaining boundary}

The accompanying Python program constructs the Pochhammer spectral
coefficients and the direct logarithmic coefficients independently over
the rational numbers. It compares the resulting ordinary remainders
with independent finite-head Hurwitz continuations, evaluates the
resolvent quadrature directly, checks the nonzero third-order correction,
and compares the convergent confluent series with its integral.
Numerical diagnostics are distinct from the all-order proofs above;
they are not interval certificates.

The answered repository question is the several-mass second derivative
and its confluent limit, with the ordinary logarithmic cross term made
explicit. The one-dimensional integral in \eqref{msj:eq:quadrature}
is itself an exact evaluation coordinate. A finite reduction of that
integral to a specified algebra of ordinary polylogarithms for arbitrary
$a,B,C$ remains a separate question. For odd or noninteger spectral
weights, resonant odd powers can occur and additional local Laurent
terms must be retained; the parity-protected bilateral formula above
does not silently assert their absence.

% END EMBEDDED sections/05_bilateral.tex


% BEGIN EMBEDDED sections/04_stieltjes_jets.tex
\section{Every spectral derivative of a product of quadratic factors}
\label{sec:onesided}

The preceding quadratic-spectrum report asks both for higher factorization
derivatives and for the logarithmic cross term for several masses
\cite[Section 7, questions 2 and 7]{Quadratic}.  We give a finite
subtraction formula at every derivative order, and then state its
conjugate-linear specialization to the precise defect in question 2.
Its one-sided version
also identifies exactly which generalized Stieltjes constants occur.
First-derivative multiplicative anomalies are classical
\cite{EVZ,PolynomialAnomaly}; the result here concerns the specified
higher derivatives and their explicit normalization.

Fix $a>0$, an integer $r\ge2$, nonnegative real numbers
$B_1,\ldots,B_r$, and a polynomial
$W(x)=\sum_{d=0}^{D}w_dx^d$.  Write $x_n=n+a$ and
$Q_j(x)=x^2+B_j$.  In their initial half-planes of absolute convergence set
\begin{align}
 E_{W,j}(s;a)&=\sum_{n\ge0}W(x_n)Q_j(x_n)^{-s},\\
 Z_{W,\boldsymbol B}(s;a)&=
       \sum_{n\ge0}W(x_n)\prod_{j=1}^rQ_j(x_n)^{-s}.
\end{align}
For the quadratic factors $Q_j$, logarithms are real logarithms of
positive numbers.  The conjugate-linear specialization below explicitly
uses principal complex logarithms.
The normalized difference is
\begin{equation}\label{one:defect}
 \mathfrak D_W(s)=\frac1r\sum_{j=1}^rE_{W,j}(rs;a)
                         -Z_{W,\boldsymbol B}(s;a),\qquad
 \Delta_m(W)=\mathfrak D_W^{(m)}(0).
\end{equation}
Thus $\Delta_m=r^{m-1}\sum_jE_{W,j}^{(m)}(0)-Z_{W,\boldsymbol B}^{(m)}(0)$.
The factor $r^{m-1}$ is part of the definition, and cannot be omitted
when comparing different total spectral degrees.

\subsection{A finite coefficient algebra}
Define polynomials $h_k(s)$ by the Taylor expansion at $t=0$,
\begin{equation}\label{one:h}
 h(s,t)=\frac1r\sum_{j=1}^r(1+B_jt)^{-rs}
                -\prod_{j=1}^r(1+B_jt)^{-s}
       =\sum_{k\ge2}h_k(s)t^k.
\end{equation}
An entirely finite formula for each coefficient is
\begin{equation}\label{one:hk}
 h_k(s)=(-1)^k\left\{
 \frac{(rs)_k}{r\,k!}\sum_jB_j^k
 -\sum_{k_1+\cdots+k_r=k}
                \prod_{j=1}^r\frac{(s)_{k_j}B_j^{k_j}}{k_j!}
                  \right\}.
\end{equation}
Here $(s)_0=1$.  Put $A_{k,j}=h_k^{(j)}(0)$; it is zero for $j>k$.
Both $h_0$ and $h_1$ vanish identically, and every $h_k$ is divisible
by $s^2$.  The latter assertion follows before coefficient extraction:
$h(0,t)=\partial_sh(0,t)=0$.  In particular, if
\begin{equation}\label{one:variance}
 V(\boldsymbol B)=\sum_{i<j}(B_i-B_j)^2,
\end{equation}
then $h_2(s)=V(\boldsymbol B)s^2/2$.

\begin{theorem}[Finite subtraction at all spectral orders]
\label{thm:alljets}
Choose any integer
\begin{equation}\label{one:K}
 K\ge\max\left(1,\left\lfloor\frac{D+1}{2}\right\rfloor\right).
\end{equation}
The functions in \eqref{one:defect} are holomorphic at zero, and
\begin{align}\label{one:master}
 \mathfrak D_W(s)={}&\mathcal R_K(s)
   +\sum_{d=0}^{D}w_d\sum_{k=2}^{K}
                     h_k(s)\zeta(2rs+2k-d,a),\\
 \mathcal R_K(s)={}&\sum_{n\ge0}W(x_n)x_n^{-2rs}
       \left\{h(s,x_n^{-2})-\sum_{k=2}^{K}h_k(s)x_n^{-2k}\right\}.
       \label{one:remainder}
\end{align}
The remainder is holomorphic in
$\Re s>(D-2K-1)/(2r)$, with normally convergent fixed spectral
derivatives there.  Every apparent pole of the finite sum at $s=0$
is removable after multiplication by $h_k(s)$.

For $m\ge1$ define
\begin{align}
 L_j(x)&=\log(x^2+B_j),\\
 T_m(x)&=(-1)^m\left\{r^{m-1}\sum_j L_j(x)^m
                         -\left(\sum_jL_j(x)\right)^m\right\},
                         \label{one:T}\\
 U_{k,m}(X)&=\sum_{j=2}^{m}\binom mj A_{k,j}(-2rX)^{m-j}.
                         \label{one:U}
\end{align}
An empty sum is zero.  Then the following is an ordinary convergent
series plus a finite, explicit correction:
\begin{equation}\label{one:jet}
 \boxed{\begin{aligned}
 \Delta_m(W)={}&\sum_{n\ge0}W(x_n)
       \left\{T_m(x_n)-\sum_{k=2}^{K}x_n^{-2k}
                                      U_{k,m}(\log x_n)\right\}\\
 &+\sum_{d=0}^{D}w_d\sum_{k=2}^{K}\mathcal C_{m,k}(2k-d;a),
 \end{aligned}}
\end{equation}
where, if $\nu\ne1$,
\begin{equation}\label{one:regularC}
 \mathcal C_{m,k}(\nu;a)=
       \sum_{j=2}^{m}\binom mj A_{k,j}(2r)^{m-j}
                                     \zeta^{(m-j)}(\nu,a),
\end{equation}
and at the resonant value $\nu=1$,
\begin{equation}\label{one:resonantC}
 \boxed{\displaystyle
 \mathcal C_{m,k}(1;a)=\frac{A_{k,m+1}}{2r(m+1)}
       +\sum_{j=2}^{m}\binom mj A_{k,j}(-2r)^{m-j}
                                     \gamma_{m-j}(a).}
\end{equation}
Consequently only $\gamma_0(a),\ldots,\gamma_{m-2}(a)$ can occur at
the $m$th derivative.  They occur only through odd weight monomials
$x^{2k-1}$ with $k\ge2$.
\end{theorem}

\begin{proof}
For sufficiently large $\Re s$, subtract the finite Taylor polynomial
in \eqref{one:h} from each summand in \eqref{one:defect}; the removed
powers sum to the Hurwitz functions in \eqref{one:master}.  This proves
the equality where all series converge absolutely.

On a compact set of $s$ and masses, the Taylor remainder is
$O(x^{-2K-2})$ as $x\to\infty$, uniformly with every fixed spectral
derivative.  A derivative of $x^{-2rs}$ introduces only a power of
$\log x$.  Hence the summands of \eqref{one:remainder} are bounded
by $C x^{D-2K-2-2r\Re s}(1+\log x)^M$, which is summable in the
stated half-plane.  This argument uses the Taylor expansion only for
large $x$.  The finitely many small $x_n$ are kept in the original
expression, so no assumption such as $B_j<a^2$ is needed.
It proves normal convergence, holomorphy, and the legitimacy of
termwise differentiation.  It also continues the original identity
from its initial half-plane to a neighborhood of zero.

The same finite-subtraction argument applied separately to
$\prod_j(1+B_jt)^{-s}$ and $(1+B_jt)^{-s}$ proves that the original
spectral functions are holomorphic at zero.  Their positive-$t$-degree
coefficients vanish at $s=0$ and thus cancel any simple Hurwitz pole;
their degree-zero terms are $\zeta(-d,a)$, regular at zero.

The only finite term in \eqref{one:master} with a pole at zero has
$2k-d=1$.  The Stieltjes convention is
\begin{equation}\label{one:stieltjes}
 \zeta(1+u,a)=\frac1u+
                 \sum_{j\ge0}\frac{(-1)^j}{j!}\gamma_j(a)u^j.
\end{equation}
Since $h_k(s)=O(s^2)$, the product is holomorphic.  At a nonresonant
integer, the ordinary product rule gives \eqref{one:regularC}.
At a resonance, the pole is $1/(2rs)$, so its $m$th coefficient
contributes $A_{k,m+1}/(2r(m+1))$.  The regular coefficients in
\eqref{one:stieltjes} give the remaining terms of
\eqref{one:resonantC}.  This proves the finite correction exactly,
including its sign and the extra local term.

Finally, differentiating the actual summands in the remainder gives
\eqref{one:T} and \eqref{one:U}.  The normal convergence just proved
allows the differentiation, and yields \eqref{one:jet}.  Increasing
$K$ merely transfers an exactly summed term between the remainder
and the correction, so the answer is independent of that choice.
\end{proof}

\subsection{The local first derivative and the nonlocal second derivative}

The first derivative is particularly short:
\begin{corollary}[First-derivative residue]\label{cor:firstanomaly}
For the convention \eqref{one:defect},
\begin{equation}\label{one:first}
 \Delta_1(W)=\frac1{4r}[x^{-1}]\,W(x)
       \sum_{i<j}\left\{\log\frac{x^2+B_i}{x^2+B_j}\right\}^{\!2}.
\end{equation}
It is independent of $a$.  If $\deg W\le2$, it vanishes for all masses.
For $W=x^3$, it equals $V(\boldsymbol B)/(4r)$.
\end{corollary}
\begin{proof}
The remainder and nonresonant corrections start at order $s^2$.
Moreover $\partial_s^2h(0,t)$ equals the sum of the pairwise squared
logarithmic differences in \eqref{one:first}.  Substitute $m=1$
in \eqref{one:resonantC}.
\end{proof}

At the second derivative the pairwise logarithmic squares are still
present, but are no longer purely local.  In addition, at resonances
one must retain $A_{k,3}/(6r)$.  The next examples make both features
explicit.

\begin{corollary}[A single Stieltjes coefficient at every derivative]
\label{cor:singleStieltjes}
Take $r=2$, $W=x^3$, $B\ne C$, and set $\delta=B-C$.
For every integer $m\ge2$, put
\begin{equation}\label{one:cm}
 c_m=\binom m2(-4)^{m-2}\delta^2.
\end{equation}
With $T_m$ as in \eqref{one:T},
\begin{equation}\label{one:isolated}
 \boxed{\displaystyle
 \Delta_m(x^3)=
   \sum_{n\ge0}\left\{x_n^3T_m(x_n)
                -c_m\frac{\log^{m-2}x_n}{x_n}\right\}
       +c_m\gamma_{m-2}(a).}
\end{equation}
The displayed series converges absolutely and locally uniformly in
$a>0$ and the masses.  Also $\Delta_1(x^3)=\delta^2/8$.
\end{corollary}
\begin{proof}
Choose $K=2$.  Here $h_2(s)=\delta^2s^2/2$ exactly, so
$A_{2,2}=\delta^2$ and $A_{2,j}=0$ for $j>2$.
There is just the single resonant integer $2k-d=1$.  Equations
\eqref{one:jet} and \eqref{one:resonantC} give the result.
\end{proof}

For clarity the first two instances are
\begin{align}
 \Delta_2(x^3)&=\sum_{n\ge0}\left\{
      x_n^3\log^2\frac{x_n^2+B}{x_n^2+C}-\frac{\delta^2}{x_n}\right\}
                         -\delta^2\psi(a),\label{one:x3second}\\
 \Delta_3(x^3)&=\sum_{n\ge0}\left\{
      -3x_n^3\log((x_n^2+B)(x_n^2+C))
                     \log^2\frac{x_n^2+B}{x_n^2+C}
                 +12\delta^2\frac{\log x_n}{x_n}\right\}
                         -12\delta^2\gamma_1(a).\label{one:x3third}
\end{align}
In \eqref{one:x3second} the summand is $O(x_n^{-3})$; in
\eqref{one:x3third} it is $O(x_n^{-3}\log x_n)$.
The all-order formula isolates one prescribed generalized Stieltjes
constant, instead of introducing an undefined finite part of a
logarithmic series.

There is already an additional local correction at the next odd weight.
For $W=x^5$ and $r=2$,
\begin{equation}\label{one:x5second}
 \begin{split}
 \Delta_2(x^5)={}&\sum_{n\ge0}\left\{
 x_n^5\log^2\frac{x_n^2+B}{x_n^2+C}
       -\delta^2x_n+\frac{\delta^2(B+C)}{x_n}\right\}\\
 &+\delta^2\zeta(-1,a)
       -\delta^2(B+C)\gamma_0(a)-\frac{\delta^2(B+C)}4.
 \end{split}
\end{equation}
Indeed $A_{3,2}=-\delta^2(B+C)$ and
$A_{3,3}=-3\delta^2(B+C)$.  The final term in
\eqref{one:x5second} comes from the latter coefficient multiplied
by the Hurwitz pole.  Simply assigning finite parts to the squared
logarithm would miss it.

\subsection{The source's conjugate-linear factorization problem}
\label{one:linear-section}

The second question in \cite{Quadratic} specifically factors a single
quadratic into its two conjugate linear factors.  This requires a
different spectral scale from a product of two quadratics.  Here is the
explicit specialization, including that scale.  For $a,b>0$ define
\[
 F_{W,\pm}(s)=\sum_{n\ge0}W(x_n)(x_n\pm ib)^{-s},\qquad
 Z_W(s)=\sum_{n\ge0}W(x_n)(x_n^2+b^2)^{-s},
\]
using principal logarithms.  Since $x_n>0$,
$\log(x_n+ib)+\log(x_n-ib)=\log(x_n^2+b^2)$ without a branch term.
Put
\[
 \widehat{\mathfrak D}_W(s)=\tfrac12
        \{F_{W,+}(2s)+F_{W,-}(2s)\}-Z_W(s),\qquad
 \widehat\Delta_m=\widehat{\mathfrak D}_W^{(m)}(0).
\]
The relevant coefficient polynomials are
\begin{align}
 \widehat h(s,t)
 &=\tfrac12\{(1+ibt)^{-2s}+(1-ibt)^{-2s}\}
                              -(1+b^2t^2)^{-s}
    =\sum_{k\ge1}\eta_k(s)t^{2k},\label{one:linear-h}\\
 \eta_k(s)&=(-1)^kb^{2k}
       \left\{\frac{(2s)_{2k}}{(2k)!}-\frac{(s)_k}{k!}\right\},
 \qquad \widehat A_{k,j}=\eta_k^{(j)}(0).
 \label{one:linear-eta}
\end{align}
Every $\eta_k$ is divisible by $s^2$, although its degree can now
be $2k$.  In particular $\eta_1(s)=-2b^2s^2$.

\begin{corollary}[Every conjugate-linear spectral defect]
\label{cor:linear-jets}
Choose $K$ as in \eqref{one:K}.  The complete formula
\eqref{one:jet} holds for $\widehat\Delta_m$ with these replacements:
the finite coefficient sums begin at $k=1$; $A_{k,j}$ is replaced by
$\widehat A_{k,j}$; every spectral scale $2r$ is replaced by $2$;
and the raw derivative becomes
\begin{equation}\label{one:linear-T}
 \widehat T_m(x)=(-1)^m\left\{
   2^{m-1}(\ell_+(x)^m+\ell_-(x)^m)
                    -(\ell_+(x)+\ell_-(x))^m\right\},\qquad
 \ell_\pm(x)=\log(x\pm ib).
\end{equation}
In particular the resonant correction is exactly
\begin{equation}\label{one:linear-contact}
 \widehat{\mathcal C}_{m,k}(1;a)
 =\frac{\widehat A_{k,m+1}}{2(m+1)}
    +\sum_{j=2}^m\binom mj\widehat A_{k,j}(-2)^{m-j}
                                      \gamma_{m-j}(a).
\end{equation}
All remainder sums are ordinary and absolutely convergent.
\end{corollary}
\begin{proof}
The exact finite-subtraction identity is
\begin{align*}
 \widehat{\mathfrak D}_W(s)
 ={}&\sum_{n\ge0}W(x_n)x_n^{-2s}
       \left\{\widehat h(s,x_n^{-1})-
                         \sum_{k=1}^K\eta_k(s)x_n^{-2k}\right\}\\
 &+\sum_{d=0}^D w_d\sum_{k=1}^K
                            \eta_k(s)\zeta(2s+2k-d,a).
\end{align*}
Its tail is normally convergent for $\Re s>(D-2K-1)/2$.
The finitely many terms with $x_n\le b$ are kept exactly, so the
binomial expansion is used only within its tail disk.  The principal
logarithms are analytic because the linear factors stay in the right
half-plane.  Taylor differentiation of this identity is the proof of
Theorem~\ref{thm:alljets} with scale $2$ and coefficient
\eqref{one:linear-eta}; its pole is $1/(2s)$.
This proves the stated correction and convergence claims.
\end{proof}

For the first odd weight $W=x$, put
$\widehat c_m=-4b^2\binom m2(-2)^{m-2}$ for $m\ge2$.
Then the exact all-order identity simplifies to
\begin{equation}\label{one:linear-isolate}
 \widehat\Delta_m(x)=\sum_{n\ge0}\left\{
 x_n\widehat T_m(x_n)-\widehat c_m
                         \frac{\log^{m-2}x_n}{x_n}\right\}
                         +\widehat c_m\gamma_{m-2}(a),
 \qquad \widehat\Delta_1(x)=-b^2.
\end{equation}
At order two, $\widehat T_2(x)=-4\arctan^2(b/x)$, so
\begin{equation}\label{one:linear-second}
 \boxed{\displaystyle
 \widehat\Delta_2(x)=
 4\sum_{n\ge0}\left\{\frac{b^2}{x_n}
                       -x_n\arctan^2\frac b{x_n}\right\}
                                      +4b^2\psi(a).}
\end{equation}
Evenness of $W$ removes the first-derivative local anomaly but does
not remove the higher nonlocal part.  Already at $W=1$,
\[
 \widehat\Delta_2(1)=-4\sum_{n\ge0}\arctan^2\frac b{n+a}.
\]
The individual linear terms themselves have finite Hurwitz evaluations:
\begin{equation}\label{one:linear-hurwitz}
 F_{W,\pm}^{(m)}(0)=
 \sum_{d=0}^D w_d\sum_{j=0}^d\binom dj
       (\mp ib)^{d-j}\zeta^{(m)}(-j,a\pm ib).
\end{equation}
Thus the corollary evaluates every $Z_W^{(m)}(0)$ using a finite
complex-Hurwitz part and the explicitly convergent remainder.
For example,
\begin{equation}\label{one:linear-unweighted}
 Z_1''(0)=4\Re\zeta''(0,a+ib)
                     +4\sum_{n\ge0}\arctan^2\frac b{n+a}.
\end{equation}
Equations \eqref{one:linear-contact} and \eqref{one:linear-T}
answer the stated higher-jet factorization problem without assuming that
its first-derivative parity criterion persists.

\subsection{A digamma integral and an anchored dilogarithm primitive}

The cubic-weight second derivative has a compact ordinary integral.
For $B>C\ge0$,
\begin{equation}\label{one:digammaintegral}
 \boxed{\displaystyle
 \Delta_2(x^3)=2\int_C^B
       t\,\Re\psi(a+i\sqrt t)
                       \log\frac{B-t}{t-C}\,dt.}
\end{equation}
The endpoint logarithms are integrable.  To prove it, first use
\[
 \log^2\frac{x^2+B}{x^2+C}
       =\int_C^B\!\int_C^B\frac{ds\,dt}{(x^2+s)(x^2+t)}.
\]
The normally convergent digamma difference is
\[
 \sum_{n\ge0}\left\{\frac{x_n}{x_n^2+t}-\frac1{x_n}\right\}
        =\psi(a)-\Re\psi(a+i\sqrt t).
\]
Partial fractions first give the subtracted double sum
$\psi(a)+(f(s)-f(t))/(t-s)$, where
$f(t)=t\Re\psi(a+i\sqrt t)$.  The term $\delta^2\psi(a)$
from integrating its constant part is cancelled by the explicit
$-\delta^2\psi(a)$ in \eqref{one:x3second}.  The remaining kernel is
$(f(s)-f(t))/(t-s)$.
The diagonal is removable.  A symmetric exclusion of $s=t$, followed
by elementary integration of $1/(t-s)$, gives
$2\int_C^B f(t)\log((B-t)/(t-C))dt$.  This proves
\eqref{one:digammaintegral} without separating divergent series.

For an explicit primitive in the shift, set $\delta=B-C>0$,
$y=(x^2+B)/(x^2+C)>1$, and define
\begin{align}
 J_2(y)&=-\frac{y\log^2y}{y-1}-2\Li_2(1-y),\label{one:J2}\\
 J_3(y)&=-\frac{\log^2y}{2(y-1)^2}+\frac12\log^2y
       +\Li_2(1-y)-\frac{\log y}{y-1}
       +\log(y-1)-\log y,\label{one:J3}\\
 \Phi_{B,C}(x)&=-\frac{\delta^2}{2}J_3(y)
       +\frac{C\delta}{2}J_2(y)-\delta^2\log x
       +\frac{\delta^2}{2}\log\delta-\frac34\delta^2.
       \label{one:Phi}
\end{align}
Here $1-y<0$, so the dilogarithm is real and unambiguous.

\begin{theorem}[Convergent dilogarithm--Gamma primitive]
\label{thm:primitive}
For $a,a_0>0$ and $B>C\ge0$,
\begin{equation}\label{one:primitive}
 \boxed{\begin{aligned}
 \int_{a_0}^{a}\Delta_2(x^3;u,B,C)\,du
 ={}&\sum_{n\ge0}\{\Phi_{B,C}(n+a)-\Phi_{B,C}(n+a_0)\}\\
 &-\delta^2\log\frac{\Gamma(a)}{\Gamma(a_0)}.
 \end{aligned}}
\end{equation}
Even the two sums of $\Phi$ separately converge absolutely.  Every
normalization constant has been fixed in \eqref{one:Phi}.
\end{theorem}
\begin{proof}
Direct differentiation using
$\Li_2'(z)=-\log(1-z)/z$ gives
$J_2'(y)=\log^2y/(y-1)^2$ and
$J_3'(y)=\log^2y/(y-1)^3$.
The substitution $v=x^2=\delta/(y-1)-C$ consequently proves
\[
 \Phi_{B,C}'(x)=x^3\log^2\frac{x^2+B}{x^2+C}-\frac{\delta^2}{x}.
\]
As $y=1+\varepsilon\downarrow1$, $J_2(y)\to0$ and
$J_3(y)=\log\varepsilon-3/2+O(\varepsilon)$.
These expansions show that the last two constants in
\eqref{one:Phi} make $\Phi_{B,C}(x)\to0$ as $x\to\infty$.
Since its derivative is $O(x^{-3})$, one has
$\Phi_{B,C}(x)=O(x^{-2})$.  Its series and its differentiated
series are therefore locally normally convergent for $a>0$.
Sum the derivative and use \eqref{one:x3second} together with
$(\log\Gamma)'=\psi$.  Evaluation at $a_0$ fixes the anchor and
proves the displayed identity.
\end{proof}

The higher Stieltjes terms in \eqref{one:isolated} also have the
classical anchored primitives
\begin{equation}\label{one:Stieltjesprimitive}
 \int_{a_0}^a\gamma_n(u)\,du
   =\frac{(-1)^{n+1}}{n+1}
       \{\zeta^{(n+1)}(0,a)-\zeta^{(n+1)}(0,a_0)\}.
\end{equation}
This follows by expanding
$\partial_a\zeta(s,a)=-s\zeta(s+1,a)$ at $s=0$; it is not a
new identity claimed by this article.  Equation \eqref{one:primitive}
provides the additional explicit dilogarithmic primitive for the
convergent nonlinear part at the second derivative.

% END EMBEDDED sections/04_stieltjes_jets.tex


% BEGIN EMBEDDED sections/07_audit.tex
\section{Source audit and proposed corrections}
\label{sec:audit}

The inspected sources support the preceding developments, but the
research-status map needs an additional cross-reference.  The correction
below concerns attribution and an already answered construction problem.
It does not invalidate the same-direction spectral representations.
No counterexample to the inspected source reconstruction or quadratic
factorization theorems was found.

\subsection{The connected double family is an established interpolant}

For $\Re q>0$ write
\[
 Z_q(s_1,\ldots,s_d)=
  \sum_{n_1>\cdots>n_d\ge0}\prod_{j=1}^d(n_j+q)^{-s_j},
 \qquad Z_q(\varnothing)=1,
\]
initially in an absolute-convergence domain.  The same-direction report
defines
\begin{equation}\label{audit:connected}
 \mathcal C_{A,B}(q)=\zeta(A,B)-Z_q(A,B)
                      +\zeta(A,q)\{\zeta(B,q)-\zeta(B)\}.
\end{equation}
The earlier independent-orders report defines the relative interpolant
by ordered convolution \cite{IndependentOrders}:
\begin{equation}\label{audit:relative}
 Z_b(w)=\sum_{w=uv}Z_a(u)\mathcal H_{a,b}(v),
 \qquad \mathcal H_{a,b}(\varnothing)=1.
\end{equation}
The sum is over prefix--suffix decompositions of the word.  At depth two,
\begin{align*}
 \mathcal H_{q,1}(A,B)
 &=\zeta(A,B)-Z_q(A,B)
                  -\zeta(A,q)\{\zeta(B)-\zeta(B,q)\}.
\end{align*}
Thus there is an exact identity of the full functions,
\begin{equation}\label{audit:identification}
 \boxed{\mathcal C_{A,B}(q)=\mathcal H_{q,1}(A,B)
                              =\Psi(A,B;q-1).}
\end{equation}
The final member is the Ihara--Nakamura--Yamamoto interpolant
\cite[Proposition 2.1 and Corollary 4.2]{INY}.  Their theorem already
gives entire dependence on all independent indices at arbitrary depth.
The relative report explicitly cites that result and develops its
two-endpoint form.  This identification is algebraic on an open
convergence domain, so it does not rely on uniqueness from values at
integer $q$.

For clarity, the established general finite expression is
\begin{equation}\label{audit:star}
 \mathcal H_{q,1}(s_1,\ldots,s_d)=
  \sum_{j=0}^d(-1)^j
   Z_q^\star(s_j,\ldots,s_1)\,
                         \zeta(s_{j+1},\ldots,s_d).
\end{equation}
The empty prefix and suffix have value one, and $Z_q^\star$ uses weak
inequalities.  This is the convolution inverse of the strict ordered
series.  It preserves the published entire normalization, rather than
choosing separate finite parts for its meromorphic constituents.

The corresponding divided-difference kernel is also already available
\cite{INY,HarmonicGaps}.  If $X_0=0$, $X_j=v_1+\cdots+v_j$, then an
equivalent convenient form is
\begin{equation}\label{audit:kernel}
 K_q^{(d)}(v_1,\ldots,v_d)
   =(-1)^d[e^{X_0},e^{X_1},\ldots,e^{X_d}]
                         (y\mapsto y^{d-q}).
\end{equation}
Here the brackets are the divided difference in $y$, extended
confluently when nodes coincide.  Its normalized Mellin transform is
$\mathcal H_{q,1}$, initially when every $\Re s_j>0$.
At $d=2$ it is precisely the kernel of the same-direction report.
To compare with the earlier formula using nodes $e^{-X_j}$, invert
each node in its finite barycentric sum: the product of all nodes
cancels the exponential prefactor, leaving the sign $(-1)^d$ and
the power $d-q$ in \eqref{audit:kernel}.

The shift law and the depth-two stuffle law in this normalization are
\begin{align}
 \mathcal H_{q+1,1}(s_1,\ldots,s_d)
 &=\mathcal H_{q,1}(s_1,\ldots,s_d)
                   +q^{-s_1}\mathcal H_{q,1}(s_2,\ldots,s_d),
                   \label{audit:shift}\\
 h_q(s)h_q(t)
 &=\mathcal H_{q,1}(s,t)+\mathcal H_{q,1}(t,s)+h_q(s+t),
 \quad h_q(s)=\zeta(s)-\zeta(s,q).
 \label{audit:stuffle}
\end{align}
They are inherited identities for entire functions.  Sign conventions
can be checked immediately: the same-direction report uses the
depth-one difference $\zeta(s,q)-\zeta(s)=-h_q(s)$, which accounts
for the apparent change of sign in its displayed laws.

\paragraph{Proposed source edit.}
The attribution paragraph of \cite{SameDirection} should identify its
entire connected double family as the depth-two specialization of
\cite{INY,IndependentOrders}.  Its second further-research question
should no longer ask for existence, entireness, shift/stuffle laws, or
the basic higher divided-difference kernel.  Those constructions are
already available.  A useful replacement asks for explicit minimal
coordinate reductions on specified mixed nonpositive-index loci,
or for comparisons between genuinely different normalizations.
The Barnes and beta--Newton same-direction representations remain
separate results and are not affected by this correction.

\subsection{Updates justified by the new proofs}

The order-eight statement in \cite{Reconstruction} correctly retained
the congruence qualification $b\not\equiv1\pmod{16}$.  It was not a
false theorem.  After the present proof and certificate are accepted,
the collective statement may be strengthened to every $b\ge5$ and its
Q1 closed.  The original congruence proof remains useful as the finite
range reduction.  Its previously rejected simple valuation extrapolation
should remain rejected; the new proof does not restore that formula.

The same report's Q5 is now answered for the fourth power and its
regular even moments.  The statement should refer to the six displayed
formal coordinates and the all-index Bernoulli transform.  It should
not assert independence of the evaluated constants or promote fifth
and higher powers to the same status.  The starting value $V_0=Q_7$
must retain Hoffman's attribution.

For \cite{Quadratic}, Q7 is answered by the pairwise logarithmic
correction and the confluent resolvent formulas, while Q2 is answered
by Corollary~\ref{cor:linear-jets}.  The latter uses the source's
linear factors and spectral scale two.  A product of two quadratic
factors has scale four; confusing them changes every higher correction.
The explicit primitive in Theorem~\ref{thm:primitive} is progress on
the mixed primitive direction, not a classification of all shift--mass
integrals.

\subsection{Corrections that must accompany higher-order reuse}

Three points are essential when extending the existing formulas.
First, the identity $\mathcal H_a(-m)=0$ does not imply
$\mathcal H_a'(-m)=0$.  At bilateral third order the omitted term
would be exactly $-12(B-C)^2\log(2\sin\pi a)$ for $P(v)=v^2$.
Second, multiplication by a Hurwitz pole uses one more coefficient
of the vanishing factor than multiplication by its regular part.
The term $A_{k,m+1}/(2r(m+1))$ in
\eqref{one:resonantC} is therefore required.  Its first additional
second-order example is $-\delta^2(B+C)/4$ in
\eqref{one:x5second}.  Third, an antiderivative of a summand is not
automatically summable: the constant in $\Phi_{B,C}$ is fixed so
that it decays to zero.  Without it, the separate series of primitives
would diverge.

These are requirements of the new extensions, not errors falsely
attributed to source theorems that only treated first derivatives.
The article and integration notes retain this distinction.

\subsection{Unchanged arithmetic conjectures}

The canonical Gaussian $S_6$ identity and the revised $S_8$ candidate
remain conjectural in this continuation \cite{Manuscript}.  The rejected
older $S_8$ vector remains different from the revised candidate.  No
failed relation search, small residual, or finite formal rank calculation
is converted here into an arithmetic proof or a nonreduction assertion.
Likewise, no new claim is made about the transcendence or numerical
independence of the six quartic coordinates.  These boundaries should
remain visible when the results are integrated into the collective text.

% END EMBEDDED sections/07_audit.tex


% BEGIN EMBEDDED sections/08_research.tex
\section{Further research questions with exact targets}
\label{sec:research}

The following problems are organized around the structures proved above.
Each asks for an identity, a finite reduction, or a proof of a specified
conjecture.  A successful result should identify the precise function
class and normalization before a computation is interpreted.

\begin{enumerate}[leftmargin=*,label=\textbf{R\arabic*.}]
\item \textbf{Complete every fixed-order reconstruction family.}
Theorem~\ref{ert:main} supplies a finite decision procedure once $a$
is fixed. For $a=5$, the exact threshold $B_5=1920$ reduces the
smaller-order-ten problem to $6\le b<1920$; the adjacent case
$b=6$ and the gap-four case $b=7$ are already proved. Complete the
remaining cases using exact arithmetic, then apply the same procedure
at $a=6$, corresponding to smaller order twelve. A structural
factorization or a congruence may eliminate most of the finite range
before direct determinants are formed.  The target is a complete
all-$b$ theorem for each family, with every exceptional index accounted
for, rather than a larger table of experiments.

\item \textbf{Remove the fixed-order qualification.}
The leading positive determinant gives eventual nonvanishing but does
not control the lower Bernoulli degrees uniformly at small separation.
Find a determinant identity, an arithmetic invariant, or a recurrence
which prevents cancellation for all $a<b$.  The existing adjacent
and gap-four cases, together with the complete small-order families,
provide exact consistency constraints.  A proof uniform in $a$ would
settle the remaining unrestricted optimal-prefix conjecture.

\item \textbf{Give closed coefficient rows for higher-order reconstruction.}
The ordinary identities \eqref{r8c:ordinary} use a completely specified
rational inverse matrix.  Seek a product formula or a fixed-size kernel
for its rows, comparable to the known smaller-order inverse construction.
It should specify every polynomial coefficient directly and retain all
subtraction constants.  A formal reference to an unevaluated inverse
or to unknown minors is a weaker endpoint than such a closed kernel.

\item \textbf{Find the fifth- and higher-power tail cancellation.}
For each integer $d\ge5$, continue
$\sum_{n\ge0}\psi'(n+1)^d x_n^{-s}$ on its regular parity lattice.
The quartic boundary identity \eqref{qtm:boundary} suggests combining
the parity components of successive powers before applying finite
summation.  Determine which lower tail powers survive, identify a fixed
coordinate space for each $d$, and construct its anchored generator.
The dimension and the required weights should be proved, not inferred
from the first few moments.

\item \textbf{Treat the resonant parity as a separate spectral family.}
For the quartic family, the points $s=-3-2k$ are actual poles with
explicit residues $h_k$.  Determine their Laurent constants and
higher spectral coefficients by a normalization compatible with
\eqref{qtm:ordinary}.  These coefficients will involve logarithmic
moments, so the regular even-moment theorem cannot simply be
differentiated at its discrete values.  An anchored analytic generator
for the Laurent family would be a natural next object.

\item \textbf{Determine the functional depth of the quartic generator.}
The present construction gives a finite depth-at-most-three expression.
Specify allowable transformations of arguments and products, then ask
whether the particular depth-three combination reduces further.
A symbolic relation must be checked by a functional derivative
identity and an anchored endpoint.  Formal symbol or shuffle rank
should not be interpreted as arithmetic irreducibility of a value.

\item \textbf{Extend exact reconstruction to mixed and shifted tails.}
Replace the identical fourth power by
$\prod_j\zeta(s_j,n+a_j)$, beginning with positive integer orders
and a common rational shift.  Determine which parity cancellations
survive and the finite replacement for
$\zeta(-2m,1/2)=0$.  Character twists should be handled only after
the shift-dependent counterterms are explicit.  The goal is ordinary
convergent identities with a specified finite arithmetic coordinate set.

\item \textbf{Reduce the spectral logarithmic quadrature on special loci.}
For rational $a$ and selected algebraic masses, seek finite reductions
of \eqref{msj:eq:quadrature} and \eqref{one:digammaintegral} into
polylogarithm order derivatives, Gamma derivatives, or a precisely
defined multiple-Gamma family.  The ordinary integral is already an
exact answer at arbitrary parameters.  A finite special-function
reduction is an additional theorem and should not be assumed from
the first-derivative factorization.

\item \textbf{Classify noninteger-weight resonance crossings.}
Replace a polynomial monomial by $x^\rho$ in the one-sided master.
As $\rho$ crosses an odd integer, an argument of the Hurwitz function
passes through one.  Construct a joint Laurent formula in $s$ and
$\rho$ which specializes to \eqref{one:resonantC} and distinguishes
different orders of taking limits.  The local coefficient at each
crossing must match the polynomial theorem exactly.

\item \textbf{Build primitive formulas at every spectral order.}
The primitive \eqref{one:primitive} treats the nonlinear second jet
with a completely fixed dilogarithmic anchor.  Higher logarithmic
powers lead naturally to multiple polylogarithms after the same
rational substitution.  Develop a finite integration rule for each
$m$ and combine it with the classical primitives of
$\gamma_{m-2}(a)$.  Also determine when differentiation in a mass
commutes with an anchored shift primitive without an additional
endpoint term.

\item \textbf{Project the Stieltjes identities to arithmetic characters.}
At $a=p/q$, apply finite character sums to the all-order correction
formulas.  Give an exact dictionary between the resulting combinations
of $\gamma_j(p/q)$, Hurwitz derivatives, and Dirichlet-$L$ derivatives.
The reflected even and odd channels should be retained separately.
The first useful deliverable would cover small denominators with
the full logarithmic normalization in $q$.

\item \textbf{Return to the Gaussian period identities with new functional input.}
The unresolved $S_6$ and revised $S_8$ identities require a relation
that survives the known restrictions of the explored Cayley relation
systems.  The higher-depth polylogarithms appearing in the quartic
generator offer additional functional expressions, but no connection
to the target period is established here.  A useful next result would
give an explicit map from one such functional identity to the exact
canonical residual, including all lower-weight and branch terms.

\item \textbf{Formalize the finite and analytic interfaces.}
The Newton triangularization, Bernoulli modular recurrence, finite
summation identities, and resonant coefficient extraction are suitable
finite proof-assistant targets.  The separate analytic obligations are
normal convergence of the subtracted families and their agreement
with the original series on a common half-plane.  A formalization
should state those interfaces explicitly.  Validated interval
evaluation of the remaining quadratures is a different project from
formal proof and can be added independently.
\end{enumerate}

% END EMBEDDED sections/08_research.tex

\appendix

% BEGIN EMBEDDED sections/09_reproducibility.tex
\section{Reproduction, certificates, and evidence boundaries}
\label{sec:reproduce}

The package contains the modular article source, a standalone combined
\TeX{} source, the compiled PDF, exact coefficient data, the inherited
order-eight certificate with provenance, and independent Python
verification programs.  All mathematical verification programs operate
without a network connection.  They require Python 3, SymPy, and mpmath;
the article builds with a standard modern \LaTeX{} installation.

\subsection{What the finite certificates establish}

The order-eight verifier reconstructs the source four-dimensional
matrix from finite differences, clears denominators, and checks its
determinant by the defining 24 permutation products.  It replays
3,852 exact polynomial coefficient comparisons in the inherited
certificate.  The new leading Hankel factor is checked both by
coefficient extraction and by its triangular structural factor.
The constants $K$, $V$, and $C$, the exact factorial comparison at
$b=512$, and all 31 remaining indices are evaluated exactly.  The
finite indices are checked by both rational Bernoulli numbers and a
separate recurrence modulo $2^{32}$.

The all-$b$ conclusion does not follow from a finite sampling rule.
It follows from the analytic domination proof, the inherited
congruence reduction, and the exhaustive finite residual set.  The
general fixed-order theorem uses its constructive factorial comparison
to prove termination; the example thresholds do not stand in for
that termination argument.

Ten additional exact normalization comparisons use $a=1,\ldots,5$
and $b=2a,2a+1$.  They independently compare the original matrix,
its shifted Hankel block, the finite Pad\'e sums, and the retained
weighted-Bernoulli matrix.  The separate threshold script computes
every coefficient norm used in the table in
Section~\ref{ert:section} by exact arithmetic.

The artifact \src{order8_identity_b17_k0.json} specifies all integer
coefficients of the ordinary convergent extraction identity for
$(p,q)=(8,34)$.  Its generator verifies the coordinate cancellation
against the original matrix, so the example does not depend on
transcribing a long printed coefficient list.

The quartic verifier reconstructs the finite-cutoff coordinates through
Faulhaber's formula, independently of the closed $D_j$ expression.
A separate exact shuffle, stuffle, and duality calculation reduces
the multiple zeta values through weight seven.  The replay checks
25 quartic coordinates, 20 centered moments, 24 boundary identities,
24 generator coefficients, the six-coordinate determinant, and the
two coefficient identities behind the rational extraction.

The bilateral spectral verifier compares rational Pochhammer
coefficients with coefficients constructed directly from powers of
elementary logarithmic series.  It checks 798 exact rational
identities, including several numbers of masses.  The one-sided
verifier adds 142 exact symbolic identities, including the general
Hurwitz Laurent product, the additional local pole term, both
dilogarithm derivatives, the primitive normalization, and the
conjugate-linear coefficients.

The columns of the first six quartic rows correspond, in order, to
$Q_7$, $Q_5$, $\zeta(2)$, $\zeta(3)$, $\zeta(4)$, and $1$:
\[
\begingroup\renewcommand{\arraystretch}{1.6}
\begin{pmatrix}
1&0&0&0&0&0\\
-\frac1{12}&\frac13&0&6&-\frac{17}{6}&0\\
\frac7{240}&-\frac16&-\frac1{10}&-\frac32&\frac{17}{30}&-\frac1{10}\\
-\frac{31}{1344}&\frac7{48}&\frac1{168}&\frac{13}{20}&\frac{17}{672}&\frac{11}{280}\\
\frac{127}{3840}&-\frac{31}{144}&\frac{299}{2160}&-\frac{2551}{5040}&-\frac{289}{540}&-\frac{17}{840}\\
-\frac{2555}{33792}&\frac{127}{256}&-\frac{3667}{6336}&\frac{45929}{67200}&\frac{497267}{253440}&\frac{1187}{277200}
\end{pmatrix}
\endgroup
\]
Their exact determinant is $7361/49896000$.

\subsection{Independent numerical diagnostics}

The numerical work tests different representations, with every working
precision and cutoff stored in its JSON record.  The quartic checks
use 120 decimal working digits and two finite-head/centered-Hurwitz-tail
choices, $(N,J)=(120,36)$ and $(240,54)$.  For $0\le r\le8$, the
largest residual from the second choice is below
$3.6\times10^{-90}$.

The bilateral calculations use 75 decimal working digits.  The
spectral continuations and logarithmic tails are constructed separately;
the elementary resolvent quadrature and confluent expression provide
additional comparisons.  In particular, at $a=1/2$ and masses $4,1$,
\[
 S_1=4.25062684984826912051041637407031330017\ldots,
\]
whereas at $P(v)=v^2$,
\[
 S_{v^2}=-66.37570183951562430903381244813937818\ldots.
\]
These values illustrate that the subtraction and the polynomial weight
are part of the identity; the sign of the unweighted logarithmic square
is not a sign rule for every subtracted weighted sum.

The one-sided diagnostics use 85 decimal working digits and 22 strict
comparisons, with tolerance $10^{-60}$.  Individual spectral families
use 11 head terms and 72 tail coefficients; the elementary logarithmic
sums use 12 head terms and 64 coefficients.  The largest recorded
strict difference is below $2\times10^{-71}$.  The choices include
$(a,B,C)=(1/3,9,0)$, for which $B/a^2=81$; thus the implementation
does not depend on an invalid global small-mass expansion.  Both
orientations of the anchored primitive are compared with a separate
Gamma quadrature.  Two further large-$x$ checks test its asymptotic
normalization.

None of these decimal diagnostics is an interval certificate or a proof
of an identity.  They test the independently derived formulas and
their implementation.  The proofs of meromorphic continuation,
normal convergence, and all-index cancellation are in the article.
No proof-assistant formalization is supplied.

\subsection{Entry points and integration}

Run \src{python verification/run_all.py} to execute the default
mathematical suites in isolated subprocesses.  The exact order-eight
source replay is included in that default.  The scripts also run
individually in their respective directories.  The recorded receipts
identify their actual runtime versions and completion status.

Build the article with
\src{latexmk -pdf -interaction=nonstopmode -halt-on-error article.tex}.
Alternatively compile \src{article_singlefile.tex} twice with
pdf\LaTeX{}.  The combined source embeds every section and bibliography
and has no external mathematical-input dependency.

The integration notes specify the source claims to update, the theorem
boundaries to retain, and suggested manuscript destinations.  Historical
incoming reports need not be overwritten: the new proof can be added
as a continuation with explicit cross-references.  The machine-readable
\src{CLAIM_STATUS.json} separates new proved claims, inherited inputs,
editorial corrections, and unresolved questions.  The final
\src{SHA256SUMS} covers all distributed source, code, data, and PDF
files other than that manifest itself.  Rebuilding may alter PDF
metadata; a new build should receive its own rendered review.

% END EMBEDDED sections/09_reproducibility.tex


% BEGIN EMBEDDED references.tex
\begin{thebibliography}{99}

\bibitem{Manuscript}
V. Reshetnikov and ProveIt Contributors,
\emph{Polylogarithms and their Arithmetic Bridges}, collective manuscript,
four volumes and associated research reports.  Inspected source revision
\sourcepin, 11 October 2026.
\url{https://github.com/VladimirReshetnikov/ProveIt/tree/73f18ef013d5b74eafc883efb2c15210642df12a/Analysis/Polylogarithms/docs/manuscript}.

\bibitem{Reconstruction}
ProveIt research collection,
\emph{Exact Reconstruction and Higher Resonances in Polygamma--Zeta Calculus}, incoming research report,
11 October 2026.  In particular Sections 2, 5, 6, and 9.
Archive \src{ProveIt_Polygamma_Reconstruction_and_Resonance_2026-10-11.zip}
in \src{docs/incoming}, at \sourcepin.
\url{https://github.com/VladimirReshetnikov/ProveIt/tree/73f18ef013d5b74eafc883efb2c15210642df12a/docs/incoming}.

\bibitem{GeneralHankel}
ProveIt research collection, Bernoulli--Tornheim--digamma continuation,
\emph{The general Hankel reduction}, source section
\src{Analysis/Polylogarithms/docs/reports/stieltjes-harmonic-resonance/continuations/bernoulli-tornheim-digamma/sections/03b_general_hankel.tex},
at \sourcepin.  The general small determinant and Pad\'e row
reduction are reproduced in Section~\ref{ert:section}.

\bibitem{Square}
ProveIt Contributors, \emph{Centered trigamma moments}, manuscript source
\src{Analysis/Polylogarithms/docs/manuscript/chapters/07-centered-trigamma-moments.tex},
at \sourcepin.  The all-even square moment formula and its
Li$_2$/Li$_3$/Li$_4$ generator.

\bibitem{Quadratic}
ProveIt research collection, \emph{Quadratic-Spectrum Stieltjes
Identities}, incoming research report, 10 October 2026.  In particular
Sections 5--7 and further-research questions 2 and 7.
Archive \src{ProveIt_Quadratic_Spectrum_Stieltjes_2026-10-10.zip}
in \src{docs/incoming}, at \sourcepin.

\bibitem{SameDirection}
ProveIt research collection, \emph{Same-Direction Polylogarithm Products},
incoming research report, 10 October 2026.  Connected coordinates,
same-direction representations, attribution paragraph, and
further-research question 2.
Archive \src{ProveIt_Same_Direction_Complex_Orders_2026-10-10.zip}
in \src{docs/incoming}, at \sourcepin.

\bibitem{IndependentOrders}
ProveIt research collection, independent-orders continuation,
\emph{An entire normalization for all independent harmonic exponents},
source section
\src{Analysis/Polylogarithms/docs/reports/stieltjes-harmonic-resonance/continuations/independent-orders/sections/02_hurwitz.tex},
at \sourcepin.  The relative interpolant, its exact identification
with the published interpolant, and the endpoint and stuffle laws.

\bibitem{HarmonicGaps}
ProveIt research collection, harmonic-gaps and mixed-moments continuation,
\emph{A jointly entire kernel for all zero-block decorations},
source section
\src{Analysis/Polylogarithms/docs/reports/stieltjes-harmonic-resonance/continuations/harmonic-gaps-mixed-moments/sections/02_joint_generator.tex},
at \sourcepin.  In particular the established divided-difference
kernel and its normalized Mellin transform.

\bibitem{Hoffman}
M. E. Hoffman, \emph{Sums of products of Riemann zeta tails},
Mediterranean Journal of Mathematics \textbf{13} (2016), 2771--2781.
\url{https://doi.org/10.1007/s00009-015-0653-9}.
Preprint \url{https://arxiv.org/abs/1602.03399}.
The fourth-power evaluation is printed immediately after Corollary 4.4.

\bibitem{INY}
K. Ihara, Y. Nakamura, and S. Yamamoto,
\emph{Interpolant of truncated multiple zeta functions},
The Ramanujan Journal \textbf{67} (2025), article 1.
Published 4 March 2025.
\url{https://doi.org/10.1007/s11139-025-01045-2}.

\bibitem{EVZ}
E. Elizalde, L. Vanzo, and S. Zerbini,
\emph{Zeta-function regularization, the multiplicative anomaly and
the Wodzicki residue}, preprint (1997).
\url{https://arxiv.org/abs/hep-th/9701060}.

\bibitem{PolynomialAnomaly}
E. G\"urel,
\emph{Resolution of multiplicative anomaly of zeta regularization for
polynomials}, arXiv:2504.14563, version 2, 3 September 2025.
\url{https://arxiv.org/abs/2504.14563}.

\bibitem{DLMF}
F. W. J. Olver et al., eds., \emph{NIST Digital Library of Mathematical
Functions}.  In particular Sections 5.5 and 5.15 (Gamma and polygamma),
24.8 (Bernoulli Fourier expansions), and 25.2, 25.11, 25.12 (zeta,
Hurwitz zeta, and polylogarithms).  Accessed 11 October 2026.
\url{https://dlmf.nist.gov/}.

\end{thebibliography}

% END EMBEDDED references.tex

\end{document}
